





\documentclass[11pt]{article}
\usepackage[T1]{fontenc}
\usepackage[utf8]{inputenc}
\usepackage{lmodern}
\usepackage[english]{babel}
\usepackage{amsmath,amssymb,mathtools,mathrsfs,array,booktabs,pdflscape,adjustbox,diagbox,graphicx}
\usepackage{geometry}
\usepackage{microtype}
\usepackage{newunicodechar}
\newunicodechar{—}{---}
\newunicodechar{“}{``}
\newunicodechar{”}{''}
\usepackage[hidelinks,unicode=true,pdfencoding=auto,bookmarksopen=true,bookmarksnumbered=false]{hyperref}
\geometry{a4paper,margin=1.45cm}
\setlength{\parindent}{1.25em}
\setlength{\parskip}{0pt}
\makeatletter
\renewcommand\footnotesize{\@setfontsize\footnotesize{7.5}{8.5}\selectfont}
\makeatother
\providecommand{\srcfn}[2]{%
  \begingroup
  \renewcommand{\thefootnote}{#1}%
  \footnote{#2}%
  \addtocounter{footnote}{-1}%
  \endgroup}
\providecommand{\srcfnmark}[1]{%
  \begingroup
  \renewcommand{\thefootnote}{#1}%
  \footnotemark%
  \addtocounter{footnote}{-1}%
  \endgroup}
\providecommand{\srcfntext}[2]{%
  \begingroup
  \renewcommand{\thefootnote}{#1}%
  \footnotetext{#2}%
  \endgroup}
\providecommand{\noethpIIrosette}{\makebox[9.44444pt][c]{\raisebox{-3.85pt}[0pt][0pt]{\includegraphics[width=14.726716981132075pt]{authority_rosette_native_supported_mask.png}}}}
\providecommand{\srcnumdisplay}[3][3.4em]{%
  \par\smallskip\noindent
  \begin{minipage}[t]{#1}\vspace*{0pt}#2\end{minipage}%
  \begin{minipage}[t]{\dimexpr\linewidth-#1\relax}\vspace*{0pt}\centering
  \(\displaystyle #3\)
  \end{minipage}\par\smallskip}
\providecommand{\qtext}[1]{``#1''}
\providecommand{\widebar}[1]{\overline{#1}}
\providecommand{\tightlist}{\setlength{\itemsep}{0pt}\setlength{\parskip}{0pt}}
\providecommand{\A}{\Delta}
\providecommand{\D}{\Delta}
\allowdisplaybreaks
\setlength{\emergencystretch}{10em}
\sloppy
\hbadness=10000
\vbadness=10000
\hfuzz=2pt
\vfuzz=10pt
\raggedbottom
\providecommand{\R}{\varrho}
\providecommand{\hata}{\widehat{a}}
\providecommand{\hatv}{\widehat{v}}
\providecommand{\modu}[1]{\;\operatorname{mod} #1}
\providecommand{\mcell}[1]{\ensuremath{\begin{gathered}#1\end{gathered}}}
\providecommand{\pair}[2]{(#1\,|\,#2)}
\providecommand{\eps}{\varepsilon}
\providecommand{\rhoidx}{\varrho}

\providecommand{\dd}{\mathrm d}
\providecommand{\Div}{\operatorname{Div}}
\providecommand{\G}{\mathfrak G}
\providecommand{\bd}{\bar{\delta}}
\providecommand{\p}{\partial}

\providecommand{\frS}{\mathfrak S}
\providecommand{\frp}{\mathfrak p}
\providecommand{\fra}{\mathfrak a}
\providecommand{\frb}{\mathfrak b}
\providecommand{\frc}{\mathfrak c}
\providecommand{\frf}{\mathfrak f}
\providecommand{\frr}{\mathfrak r}
\providecommand{\frm}{\mathfrak m}
\providecommand{\frD}{\mathfrak D}
\providecommand{\calA}{\mathcal A}
\providecommand{\calB}{\mathcal B}
\providecommand{\calN}{\mathcal N}


\providecommand{\p}{\partial}
\providecommand{\dd}{\mathrm d}
\providecommand{\modu}[1]{\;\operatorname{mod} #1}
\providecommand{\mM}{\mathfrak{M}}
\providecommand{\mN}{\mathfrak{N}}
\providecommand{\mL}{\mathfrak{L}}
\providecommand{\mG}{\mathfrak{G}}
\providecommand{\mH}{\mathfrak{H}}
\providecommand{\mH}{\mathfrak{H}}
\providecommand{\mA}{\mathfrak{A}}
\providecommand{\mB}{\mathfrak{B}}
\providecommand{\mX}{\mathfrak{X}}
\providecommand{\mY}{\mathfrak{Y}}
\providecommand{\ma}{\mathfrak{a}}
\providecommand{\mb}{\mathfrak{b}}
\providecommand{\mc}{\mathfrak{c}}
\providecommand{\mx}{\mathfrak{x}}
\providecommand{\my}{\mathfrak{y}}
\providecommand{\mz}{\mathfrak{z}}
\providecommand{\me}{\mathfrak{e}}
\providecommand{\mbarA}{\overline{\mathfrak{A}}}
\providecommand{\mbarB}{\overline{\mathfrak{B}}}
\providecommand{\mbar}[1]{\overline{#1}}
\providecommand{\mC}{\mathfrak{C}}
\providecommand{\mE}{\mathfrak{E}}
\providecommand{\mF}{\mathfrak{F}}
\providecommand{\mT}{\mathfrak{T}}
\providecommand{\mW}{\mathfrak{W}}
\providecommand{\md}{\mathfrak{d}}
\providecommand{\mf}{\mathfrak{f}}
\providecommand{\mm}{\mathfrak{m}}
\providecommand{\mn}{\mathfrak{n}}
\providecommand{\mo}{\mathfrak{o}}
\providecommand{\mpideal}{\mathfrak{p}}
\providecommand{\mq}{\mathfrak{q}}
\providecommand{\mt}{\mathfrak{t}}
\providecommand{\mv}{\mathfrak{v}}
\providecommand{\bara}{\overline{\mathfrak{a}}}

\providecommand{\mU}{\mathfrak{U}}
\providecommand{\mP}{\mathfrak{P}}
\providecommand{\mQ}{\mathfrak{Q}}
\providecommand{\mS}{\mathfrak{S}}
\providecommand{\mD}{\mathfrak{D}}
\providecommand{\mR}{\mathfrak{R}}
\providecommand{\mV}{\mathfrak{V}}
\providecommand{\ideal}[1]{\mathfrak{#1}}

\providecommand{\Amod}{\mathfrak A}
\providecommand{\Bmod}{\mathfrak B}
\providecommand{\Gmod}{\mathfrak G}
\providecommand{\Rmod}{\mathfrak R}
\providecommand{\Smod}{\mathfrak S}
\providecommand{\mideal}{\mathfrak m}
\providecommand{\nideal}{\mathfrak n}
\providecommand{\gideal}{\mathfrak g}
\providecommand{\bb}{\mathfrak b}
\providecommand{\cc}{\mathfrak c}
\providecommand{\zero}[1]{\equiv 0(#1)}
\providecommand{\Norm}{N}


\providecommand{\Amod}{\mathfrak A}
\providecommand{\Bmod}{\mathfrak B}
\providecommand{\Gmod}{\mathfrak G}
\providecommand{\Mmod}{\mathfrak M}
\providecommand{\Nmod}{\mathfrak N}
\providecommand{\Smod}{\mathfrak S}
\providecommand{\mideal}{\mathfrak m}
\providecommand{\nideal}{\mathfrak n}
\providecommand{\gideal}{\mathfrak g}
\providecommand{\hideal}{\mathfrak h}
\providecommand{\jideal}{\mathfrak j}
\providecommand{\tideal}{\mathfrak t}
\providecommand{\aideal}{\mathfrak a}
\providecommand{\bideal}{\mathfrak b}
\providecommand{\zero}[1]{\equiv 0\pmod{#1}}
\providecommand{\Norm}{N}


\providecommand{\frakS}{\mathfrak S}
\providecommand{\frakM}{\mathfrak M}
\providecommand{\frakG}{\mathfrak G}
\providecommand{\frakm}{\mathfrak m}
\providecommand{\frakn}{\mathfrak n}
\providecommand{\frakp}{\mathfrak p}
\providecommand{\frakq}{\mathfrak q}
\providecommand{\frakr}{\mathfrak r}
\providecommand{\fraka}{\mathfrak a}
\providecommand{\frakb}{\mathfrak b}
\providecommand{\frakc}{\mathfrak c}
\providecommand{\frakd}{\mathfrak d}
\providecommand{\frake}{\mathfrak e}
\providecommand{\frakg}{\mathfrak g}
\providecommand{\fraks}{\mathfrak s}
\providecommand{\frako}{\mathfrak o}
\providecommand{\frK}{\mathfrak K}
\providecommand{\barfrakm}{\overline{\mathfrak m}}

\providecommand{\frakh}{\mathfrak h}
\providecommand{\frakt}{\mathfrak t}

\providecommand{\frakR}{\mathfrak R}
\providecommand{\frakT}{\mathfrak T}
\providecommand{\frakH}{\mathfrak H}
\providecommand{\frakP}{\mathfrak P}
\providecommand{\frakQ}{\mathfrak Q}
\providecommand{\frakZ}{\mathfrak Z}
\providecommand{\frakN}{\mathfrak N}
\providecommand{\frakB}{\mathfrak B}
\providecommand{\frakF}{\mathfrak F}
\providecommand{\frakL}{\mathfrak L}
\providecommand{\frakW}{\mathfrak W}
\providecommand{\frakGg}{\mathfrak G}
\providecommand{\Endl}{\operatorname{End}}
\providecommand{\Gal}{\operatorname{Gal}}
\providecommand{\Trdeg}{\operatorname{trdeg}}

\providecommand{\frR}{\mathfrak R}
\providecommand{\frK}{\mathfrak K}
\providecommand{\frL}{\mathfrak L}
\providecommand{\frM}{\mathfrak M}
\providecommand{\frF}{\mathfrak F}
\providecommand{\frT}{\mathfrak T}
\providecommand{\frN}{\mathfrak N}
\providecommand{\frO}{\mathfrak O}
\providecommand{\frB}{\mathfrak B}
\providecommand{\frC}{\mathfrak C}
\providecommand{\frA}{\mathfrak A}

\providecommand{\mR}{\mathfrak{R}}
\providecommand{\bR}{\overline{\mathfrak{R}}}
\providecommand{\mK}{\mathfrak{K}}
\providecommand{\mL}{\mathfrak{L}}
\providecommand{\mS}{\mathfrak{S}}
\providecommand{\mT}{\mathfrak{T}}
\providecommand{\mZ}{\mathfrak{Z}}
\providecommand{\mO}{\mathfrak{O}}
\providecommand{\mo}{\mathfrak{o}}
\providecommand{\mD}{\mathfrak{D}}
\providecommand{\mE}{\mathfrak{E}}
\providecommand{\mP}{\mathfrak{P}}
\providecommand{\mf}{\mathfrak{f}}
\providecommand{\mpideal}{\mathfrak{p}}
\providecommand{\mq}{\mathfrak{q}}
\providecommand{\mt}{\mathfrak{t}}
\providecommand{\mOmega}{\mathfrak{\Omega}}
\providecommand{\bP}{\overline{P}}
\providecommand{\bq}{\overline{\mathfrak{q}}}
\providecommand{\bp}{\overline{\mathfrak{p}}}
\providecommand{\ba}{\overline{\mathfrak{a}}}
\providecommand{\bb}{\overline{\mathfrak{b}}}
\providecommand{\tr}{\operatorname{Tr}}
\providecommand{\Norm}{\operatorname{N}}


\providecommand{\rank}{\operatorname{rank}}

\providecommand{\frakA}{\mathfrak A}
\providecommand{\frakB}{\mathfrak B}
\providecommand{\frakM}{\mathfrak M}
\providecommand{\frakN}{\mathfrak N}
\providecommand{\frako}{\mathfrak o}
\providecommand{\fraka}{\mathfrak a}
\providecommand{\frakc}{\mathfrak c}
\providecommand{\frakp}{\mathfrak p}
\providecommand{\frakO}{\mathfrak O}
\providecommand{\frakP}{\mathfrak P}
\providecommand{\frakD}{\mathfrak D}
\providecommand{\frakG}{\mathfrak G}
\providecommand{\frakK}{\mathfrak K}
\providecommand{\Sp}{\operatorname{Sp}}
\providecommand{\rank}{\operatorname{rank}}
\providecommand{\Trdeg}{\operatorname{Trdeg}}
\providecommand{\charac}{\operatorname{char}}

\providecommand{\Om}{\Omega}
\providecommand{\Omfrakp}{\Omega_{\mathfrak p}}
\providecommand{\frakA}{\mathfrak A}
\providecommand{\frakB}{\mathfrak B}
\providecommand{\frakC}{\mathfrak C}
\providecommand{\frakG}{\mathfrak G}
\providecommand{\frakH}{\mathfrak H}
\providecommand{\frakP}{\mathfrak P}
\providecommand{\frakp}{\mathfrak p}
\providecommand{\frakq}{\mathfrak q}
\providecommand{\frakr}{\mathfrak r}
\providecommand{\frakS}{\mathfrak S}
\providecommand{\Cl}{\operatorname{Cl}}
\providecommand{\ind}{\operatorname{ind}}
\providecommand{\N}{\operatorname{N}}
\providecommand{\Mat}{\operatorname{M}}
\providecommand{\Gal}{\operatorname{Gal}}
\providecommand{\Norm}{\operatorname{N}}
\providecommand{\nr}{\operatorname{nr}}
\providecommand{\Tr}{\operatorname{Tr}}
\providecommand{\Br}{\operatorname{Br}}

\providecommand{\frR}{\mathfrak R}
\providecommand{\frS}{\mathfrak S}
\providecommand{\frT}{\mathfrak T}
\providecommand{\frM}{\mathfrak M}
\providecommand{\frN}{\mathfrak N}
\providecommand{\frA}{\mathfrak A}
\providecommand{\frB}{\mathfrak B}
\providecommand{\frX}{\mathfrak X}
\providecommand{\frY}{\mathfrak Y}
\providecommand{\frZ}{\mathfrak Z}
\providecommand{\frG}{\mathfrak G}
\providecommand{\frH}{\mathfrak H}
\providecommand{\frL}{\mathfrak L}
\providecommand{\frU}{\mathfrak U}
\providecommand{\frD}{\mathfrak D}
\providecommand{\frC}{\mathfrak C}
\providecommand{\frP}{\mathfrak P}
\providecommand{\Gcal}{\mathfrak G}
\providecommand{\Sp}{\operatorname{Sp}}

\providecommand{\Hh}{\mathfrak H}
\providecommand{\Ii}{\mathfrak I}
\providecommand{\Jj}{\mathfrak J}
\providecommand{\Aa}{\mathfrak A}
\providecommand{\Bb}{\mathfrak B}
\providecommand{\Cc}{\mathfrak C}
\providecommand{\Oo}{\mathfrak O}
\providecommand{\oo}{\mathfrak o}
\providecommand{\pp}{\mathfrak p}
\providecommand{\PP}{\mathfrak P}
\providecommand{\LL}{\mathfrak L}
\providecommand{\RR}{\mathfrak R}
\providecommand{\ZZ}{\mathfrak Z}
\providecommand{\ee}{\mathfrak e}

\providecommand{\Gg}{\mathfrak G}
\hypersetup{%
  pdftitle={Emmy Noether: Collected Mathematical Works in English},%
  pdfauthor={Emmy Noether},%
  pdfsubject={Complete maintained English corpus edition; source authority NOETH-DE-ED-0014},%
  pdfkeywords={Emmy Noether, English translation, collected works, algebra, invariant theory}%
}
\setcounter{tocdepth}{1}
\providecommand{\editionentry}[2]{%
  \phantomsection\addcontentsline{toc}{section}{#1}\label{#2}}
\providecommand{\editionpartentry}[2]{%
  \phantomsection\addcontentsline{toc}{subsection}{#1}\label{#2}}

\begin{document}
\begin{center}

{\Large\bfseries Ideal Theory in Ring Domains.}\par
\vspace{0.7em}
By\par
\vspace{0.25em}
Emmy Noether in Göttingen.\par
\vspace{0.9em}
\rule{4em}{0.4pt}\par
\vspace{0.9em}
{\large\bfseries Table of Contents.}
\end{center}
\noindent Introduction.\par
\noindent\hangindent=3.0em\hangafter=1 \S~1. \textbf{Ring Domain, Ideal, Finiteness Condition.}\par
\noindent\hangindent=3.0em\hangafter=1 \S~2. \textbf{Representation of an Ideal as the Least Common Multiple of Finitely Many Irreducible Ideals.}\par
\noindent\hangindent=3.0em\hangafter=1 \S~3. \textbf{Equality of the Number of Components in Two Different Decompositions into Irreducible Ideals.}\par
\noindent\hangindent=3.0em\hangafter=1 \S~4. \textbf{Primary Ideals. Uniqueness of the Associated Prime Ideals in Two Different Decompositions into Irreducible Ideals.}\par
\noindent\hangindent=3.0em\hangafter=1 \S~5. \textbf{Representation of an Ideal as the Least Common Multiple of Greatest Primary Ideals. Uniqueness of the Associated Prime Ideals.}\par
\noindent\hangindent=3.0em\hangafter=1 \S~6. \textbf{Unique Representation of an Ideal as the Least Common Multiple of Relatively Prime Irreducible Ideals.}\par
\noindent\hangindent=3.0em\hangafter=1 \S~7. \textbf{Uniqueness of the Isolated Ideals.}\par
\noindent\hangindent=3.0em\hangafter=1 \S~8. \textbf{Unique Representation of an Ideal as a Product of Coprime-Irreducible Ideals.}\par
\noindent\hangindent=3.0em\hangafter=1 \S~9. \textbf{Extension of the Investigation to Modules. Equality of the Number of Components in Decompositions into Irreducible Modules.}\par
\noindent\hangindent=3.0em\hangafter=1 \S~10. \textbf{Special Case of the Polynomial Domain.}\par
\noindent\hangindent=3.0em\hangafter=1 \S~11. \textbf{Examples from Number Theory and from the Theory of Differential Expressions.}\par
\noindent\hangindent=3.0em\hangafter=1 \S~12. \textbf{Example from Elementary-Divisor Theory.}\par

\begin{center}\textbf{Introduction.}\end{center}
The content of the present paper is the \emph{transfer of the decomposition theorems for the rational integers, respectively of the ideals in algebraic number fields, to ideals in arbitrary integral domains, and more generally ring domains}. To understand this transfer, let us first state the decomposition theorems for the rational integers in a form somewhat different from the usual formulation.

If in
\[
 a=p_1^{\varrho_1}p_2^{\varrho_2}\cdots p_\sigma^{\varrho_\sigma}=q_1q_2\cdots q_\sigma
\]
one regards the prime-power factors $q_i$ as the components of the decomposition, then these components have the following characteristic properties:

1. They are \emph{pairwise coprime}; but no $q_i$ can be represented as a product of pairwise coprime numbers, so that in this sense irreducibility holds. From the pairwise coprimeness it further follows that the product $q_1\cdots q_\sigma$ is equal to the least common multiple $[q_1\cdots q_\sigma]$.

2. Any two of the components, $q_i$ and $q_k$, are \emph{relatively prime}; that is, if $bq_i$ is divisible by $q_k$, then $b$ is divisible by $q_k$. In this sense too irreducibility holds.

3. Every $q$ is \emph{primary}; that is, if a product $b\cdot c$ is divisible by $q$, but $b$ is not divisible, then a power\footnote{If this power is always the first, then one is, as is well known, dealing with prime numbers.} of $c$ is divisible. The representation is moreover one by \emph{greatest primary components}, since the product of two distinct $q$ is no longer primary. Also with respect to decomposition into greatest primary components the $q$ are irreducible.

4. Every $q$ is \emph{irreducible} in the sense that it cannot be represented as the least common multiple of two proper divisors.

The connection of these primary numbers $q$ with the prime numbers $p$ consists in this: to each $q$ there exists one, and apart from sign only one, $p$ which is a divisor of $q$ and a power of which is divisible by $q$: the associated prime number. If $p^\varrho$ is the lowest power of this kind -- $\varrho$ the exponent of $q$ -- then here, in particular, $p^\varrho$ is equal to $q$. The uniqueness theorem may now be stated as follows:

\emph{For two different decompositions of a rational integer into the irreducible, greatest primary components $q$, the number of components, the associated prime numbers (up to sign), and the exponents coincide. Since $p^\varrho=q$, it follows from this also that the $q$ themselves coincide (up to sign).}

The indeterminacy arising from the sign is, as is well known, removed when one considers, instead of the numbers, the ideals derived from them (all numbers divisible by $a$); then the formulation holds in exactly the same way for the unique decomposition of the ideals of finite algebraic number fields into powers of prime ideals.

In what follows (\S~1), the basis will be a general ring domain satisfying only the finiteness condition that each ideal of the domain possess a finite ideal basis. Without such a finiteness condition, irreducible and prime ideals need not exist at all, as is shown by the domain of all algebraic integers, in which there is no decomposition into prime ideals.

It turns out that -- corresponding to the four characteristic properties of the components $q$ -- in general four separate decompositions exist, each arising from the preceding one by subdivision. The decomposition into coprime-irreducible ideals is a product representation; the other three decompositions are reduced (\S~2) representations as least common multiples. The connection between a primary ideal -- the irreducible ideals too are primary -- and the associated prime ideal also persists: to every primary ideal $\ideal Q$ there is uniquely determined an associated prime ideal $\ideal P$ which is a divisor of $\ideal Q$ and a power of which is divisible by $\ideal Q$. If $\ideal P^\varrho$ is the lowest such power -- $\varrho$ the exponent of $\ideal Q$ -- then here $\ideal P^\varrho$ need not coincide with $\ideal Q$. The uniqueness theorem is expressed as follows:

\emph{Decompositions 1 and 2 are unique; in two different decompositions 3 or 4 the number of components and the associated prime ideals coincide}\footnote{Presumably, beyond this, the exponents also coincide, and still more generally corresponding components are isomorphic.}\emph{; the isolated ideals occurring among the components (\S~7) are uniquely determined.}

For the proof of the decomposition theorems, the finiteness condition first yields the ``theorem on the finite chain'', first stated by Dedekind for finite number modules; from it the representation 4 of each ideal as the least common multiple of finitely many irreducible ideals is derived. By transforming the concept of reducibility of a component, one obtains from this the fundamental uniqueness theorem for decomposition 4 into irreducible ideals. By combining finitely many components at a time one rises to the remaining decompositions, whose uniqueness theorems then follow from the uniqueness theorem 4.

Finally it is shown (\S~9) that the representation by finitely many irreducible components still holds under weaker hypotheses; commutativity of the ring domain is not required, and it is enough to consider, in place of an ideal, a module with respect to the domain. In this more general case the equality of the number of components in two different decompositions still holds, whereas the notions prime and primary are tied to commutativity and to the ideal concept; by contrast, the notion coprime remains valid for ideals in noncommutative domains.

The simplest ring domain for which the four separate decompositions actually occur is the domain of all polynomials in $n$ variables with arbitrary complex coefficients. Here the individual decompositions can be interpreted, irrationally, by the behavior of algebraic configurations, and the uniqueness theorem for the associated prime ideals corresponds to the fundamental theorem of elimination theory on the unique decomposability of algebraic configurations into irreducible ones. Further examples are given by all finite integral domains made of polynomials (\S~10). But even the simple domain of all even numbers, more generally of all numbers divisible by a fixed number, already gives an example of partially separated decompositions (\S~11). An example of ideal theory in noncommutative domains is provided by elementary-divisor theory (\S~12), where unique decomposition into irreducible ideals, respectively classes, holds. These irreducible classes completely characterize the irreducible constituents of the elementary divisors, and may perhaps be viewed as their equivalent for domains in which the usual elementary-divisor theory fails.

On the existing literature the following is to be said. The decomposition into greatest primary ideals for the polynomial domain with arbitrary complex, respectively integral, coefficients was given by Lasker, and in individual points further developed by Macaulay.\footnote{E. Lasker, Zur Theorie der Moduln und Ideale. Math. Ann. 60 (1905), p. 20, Theorems VII and XIII. -- F. S. Macaulay, On the Resolution of a given Modular System into Primary Systems including some Properties of Hilbert Numbers. Math. Ann. 74 (1913), p. 66.} Both rely on elimination theory, and hence use the fact that a polynomial can be represented uniquely as a product of irreducible polynomials. In fact the decomposition theorems for ideals are independent of this prerequisite, as ideal theory in algebraic number fields suggests and as the present paper shows. The primary ideal is also defined by Lasker and Macaulay on the basis of concepts from elimination theory.

The decomposition into irreducible ideals and the decomposition into relatively-prime irreducible ideals do not seem to have been noticed in the literature even for the polynomial domain; only in Macaulay is there a remark on the uniqueness of isolated primary ideals.

The decomposition into coprime-irreducible ideals was given for the polynomial domain by Schmeidler,\footnote{W. Schmeidler, Über Moduln und Gruppen hyperkomplexer Größen. Math. Zeitschr. 3 (1919), p. 29.} using elimination theory for the proof of finiteness. There, however, the uniqueness theorem is stated only for classes of ideals, not for the ideals themselves. This latter uniqueness theorem appears in a joint paper,\footnote{E. Noether--W. Schmeidler, Moduln in nichtkommutativen Bereichen, insbesondere aus Differential- und Differenzenausdrücken. Math. Zeitschr. 8 (1920), p. 1.} where ideals in noncommutative polynomial domains are involved. Only the finite ideal basis is used there; hence theorems and methods remain valid for general ring domains, and are sharpened in the present paper with respect to equality of number (\S~11). The present investigations are a strong generalization and further development of the conceptual formations underlying those two papers. The essential feature of the two papers is the passage from representation as a least common multiple to an additive decomposition of the system of residue classes. Here, for simplicity of presentation, we remain with the least common multiple; the additive decomposition is then matched by the transformation of the notion of reducibility into a property of the complement (\S~3). Yet by considerations essentially corresponding to those of the joint paper, all the theorems stated here can also be understood as additive decomposition theorems for the system of residue classes and certain partial systems. This system of residue classes forms a ring of the same generality as the ring originally taken as basis; indeed every ring may be regarded as the system of residue classes of the ideal corresponding to the totality of identical relations among the ring elements; or also of a partial system of these relations, if the remaining relations are assumed to be satisfied already in the domain.

This remark also gives the place of Fraenkel's papers.\footnote{A. Fraenkel, Über die Teiler der Null und die Zerlegung von Ringen. J. f. M. 145 (1914), p. 139. Über gewisse Teilbereiche und Erweiterungen von Ringen. Habilitationsschrift, Leipzig, Teubner, 1916. Über einfache Erweiterungen zerlegbarer Ringe. J. f. M. 151 (1920), p. 121.} Fraenkel considers additive decompositions of rings that are subjected to such restrictive conditions (existence of regular elements, division by these, decomposability condition) that, for the corresponding ideal, the four decompositions coincide. Because of this coincidence, his finiteness condition, that the ideal should have only finitely many proper divisors -- from which he in part departs -- also means no stronger restriction than ours. Fraenkel's starting point is conditioned by the different, essentially algebraic, aims of his papers; by algebraic extension he then arrives at more general rings with less restrictive conditions.

\begin{center}
\S~1.\\[0.25em]
\textbf{Ring Domain, Ideal, Finiteness Condition.}
\end{center}

\textbf{1.} The domain $\Sigma$ taken as basis is to be a (commutative) ring in the abstract definition;\footnote{The definition is taken from Fraenkel's habilitation thesis, omitting his restrictive conditions 6, I and II; for this the commutative law of addition had to be included. Thus these are the laws defining a field with the invertibility of multiplication omitted.} that is, $\Sigma$ consists of a system of elements $a,b,c,\ldots,f,g,h,\ldots$ in which a relation satisfying the usual conditions is defined as equality; and in which by two operations (modes of composition), addition and multiplication, from any two ring elements $a$ and $b$ a third is always uniquely obtained, respectively as sum $a+b$ and as product $a\cdot b$. The ring and the otherwise wholly arbitrary operations must satisfy the following laws:
\begin{enumerate}
\item The associative law of addition: $(a+b)+c=a+(b+c)$.
\item The commutative law of addition: $a+b=b+a$.
\item The associative law of multiplication: $(a\cdot b)c=a(b\cdot c)$.
\item The commutative law of multiplication: $a\cdot b=b\cdot a$.
\item The distributive law: $a(b+c)=ab+ac$.
\item The law of unrestricted and unique subtraction. There is in $\Sigma$ a single element $x$ satisfying the equation $a+x=b$. (One denotes $x=b-a$.)
\end{enumerate}
From these properties follows the existence of zero; but a ring need not possess a unit; and the product of two elements can vanish without one of the factors vanishing. Rings for which, from the vanishing of a product, the vanishing of a factor always follows, and which in addition possess a unit, will be called proper integral domains. For the finite sum $a+a+\cdots+a$ we introduce the usual abbreviated notation $na$, where the integers $n$ are to be regarded merely as abbreviating signs, not as ring elements, and are defined recursively by $a=1a$, $na+a=(n+1)a$.

\textbf{2.} By an ideal $\ideal M$\footnote{Ideals are denoted by capital German letters. $\ideal M$ is meant to recall the example of the ideal in polynomials usually called a ``module'' or form module. Incidentally, \S\S~1--3 use only the module property and not the ideal property; compare \S~9.} in $\Sigma$ is meant a system of elements from $\Sigma$ satisfying the two conditions:
\begin{enumerate}
\item $\ideal M$ contains, together with $f$, also $a\cdot f$, where $a$ is an arbitrary element of $\Sigma$.
\item $\ideal M$ contains, together with $f$ and $g$, also the difference $f-g$; hence with $f$ also $nf$ for every integer $n$.
\end{enumerate}
If $f$ is an element of $\ideal M$, we express this as usual by $f\equiv0(\ideal M)$ and say that $f$ is divisible by $\ideal M$. If every element of $\ideal N$ is at the same time an element of $\ideal M$, and hence divisible by $\ideal M$, we say: $\ideal N$ is divisible by $\ideal M$; in symbols, $\ideal N\equiv0(\ideal M)$. $\ideal M$ is called a proper divisor of $\ideal N$ when it contains elements different from those of $\ideal N$, and hence is not conversely divisible by $\ideal N$. From $\ideal N\equiv0(\ideal M)$ and $\ideal M\equiv0(\ideal N)$ follows $\ideal N=\ideal M$.

The other familiar notions are also retained word for word. By the greatest common divisor of two ideals $\ideal A$ and $\ideal B$ -- $\ideal D=(\ideal A,\ideal B)$ -- we mean the totality of elements that can be represented in the form $a+b$, where $a$ runs through all elements of $\ideal A$ and $b$ through all elements of $\ideal B$; $\ideal D$ is again an ideal. Likewise, the greatest common divisor of infinitely many ideals -- $\ideal D=(\ideal A_1,\ideal A_2,\ldots,\ideal A_\nu,\ldots)$ -- is defined as the totality of elements $d$ representable as sums of elements from finitely many of the ideals at a time: $d=a_{i_1}+a_{i_2}+\cdots+a_{i_n}$; here too $\ideal D$ is again an ideal.

If, in particular, the ideal $\ideal M$ contains a finite number of elements $f_1,f_2,\ldots,f_\varrho$ such that
\[
 \ideal M=(f_1\ldots f_\varrho),\qquad
 f=a_1f_1+\cdots+a_\varrho f_\varrho+n_1f_1+\cdots+n_\varrho f_\varrho
\]
for every $f\equiv0(\ideal M)$, where the $a_i$ are quantities of the ring domain and the $n_i$ are integers, then $\ideal M$ is called a finite ideal; $f_1,
\ldots,f_\varrho$ an ideal basis.

In what follows we take as basis only rings $\Sigma$ satisfying the finiteness condition: \emph{Every ideal in $\Sigma$ is finite, hence possesses an ideal basis.}

\textbf{3.} From the finiteness condition there follows directly the result on which all subsequent considerations rest.

\textbf{Theorem I (theorem on the finite chain).}\footnote{First stated for number modules by Dedekind: Zahlentheorie, Suppl. XI, \S~172, Theorem VIII (4th ed.); our proof and the designation ``chain'' are taken from there. For ideals in polynomials, see Lasker, loc. cit., p. 56 (lemma). But in both cases the theorem finds only isolated application. Our applications rest throughout on the principle of choice.} \emph{Let $\ideal M,\ideal M_1,\ideal M_2,\ldots,\ideal M_\nu,\ldots$ be a countably infinite system of ideals in $\Sigma$, each of which is divisible by the following one. Then from some finite index $n$ onward all ideals are identical, $\ideal M_n=\ideal M_{n+1}=\cdots$. In other words: if $\ideal M,\ideal M_1,\ideal M_2,\ldots,\ideal M_s,
\ldots$ form a simply ordered chain of ideals such that every ideal is a proper divisor of the one immediately preceding it, then the chain breaks off after finitely many steps.}

Indeed, let $\ideal D=(\ideal M_1,\ideal M_2,\ldots,\ideal M_\nu,\ldots)$ be the greatest common divisor of the system, and let $f_1,\ldots,f_k$ be a basis of $\ideal D$, which always exists by the finiteness condition. From the divisibility assumption it follows that every element of $\ideal D$ is at the same time an element of one ideal of the chain; for from
\[
 f=g+h,\qquad g\equiv0(\ideal M_r),\qquad h\equiv0(\ideal M_s),\qquad (r\le s)
\]
one gets $g\equiv0(\ideal M_s)$ and hence $f\equiv0(\ideal M_s)$. The same holds when $f$ is a sum of several constituents. Hence there is a finite index $n$ such that
\[
 f_1\equiv0(\ideal M_n);\quad\ldots;\quad f_k\equiv0(\ideal M_n);\quad
 \ideal D=(f_1\ldots f_k)\equiv0(\ideal M_n).
\]
Since conversely $\ideal M_n\equiv0(\ideal D)$, we have $\ideal M_n=\ideal D$; and since further $\ideal M_n\equiv0(\ideal D)$ and $\ideal D=\ideal M_n\equiv0(\ideal M_r)$, we also have $\ideal M_r=\ideal D=\ideal M_n$ for every $r>n$, proving the theorem.

We note that conversely the existence of an ideal basis again follows from this theorem, so that the finiteness condition could also have been stated in this basis-free form.

\begin{center}
\S~2.\\[0.25em]
\textbf{Representation of an Ideal as the Least Common Multiple of Finitely Many Irreducible Ideals.}
\end{center}

The least common multiple $[\ideal B_1,\ideal B_2,\ldots,\ideal B_k]$ of the ideals $\ideal B_1,\ideal B_2,\ldots,\ideal B_k$ is defined as usual as the totality of the elements divisible by $\ideal B_1$, by $\ideal B_2$, $\ldots$, and by $\ideal B_k$; in symbols,
\[
 \text{from } f\equiv0(\ideal B_i),\quad (i=1,2,\ldots,k),
 \quad\text{follows}\quad
 f\equiv0([\ideal B_1,\ideal B_2,\ldots,\ideal B_k]),
\]
and conversely. The least common multiple is again an ideal; the ideals $\ideal B_i$ are also called the components of the decomposition.

\textbf{Definition I.} \emph{A representation $\ideal M=[\ideal B_1\cdots\ideal B_k]$ is called a reduced representation if no $\ideal B_i$ is absorbed into the least common multiple $\ideal A_i$ of the remaining ideals, and if no $\ideal B_i$ can be replaced by a proper divisor.\footnote{An example of a non-reduced representation is $(x^2,xy)=[(x),(x^2,xy,y^\lambda)]$ for every exponent $\lambda\ge2$; the representation $[(x),(x^2,y)]$ corresponding to $\lambda=1$ is an associated reduced one. (This representation was given to me by K. Hentzelt, who was killed in the war, as the simplest example of a nonunique decomposition into primary ideals.)} If the conditions are satisfied only for the ideal $\ideal B_i$, the representation is called reduced with respect to $\ideal B_i$. The least common multiple $\ideal A_i=[\ideal B_1,\ldots,\ideal B_{i-1},\ideal B_{i+1},\ldots,\ideal B_k]$ is called the complement of $\ideal B_i$. Representations in which only the first condition is satisfied are called shortest representations.}

For representing an ideal as a least common multiple it therefore suffices to restrict to reduced representations, by the following result.

\textbf{Lemma I.} \emph{Every representation of an ideal as the least common multiple of finitely many ideals can be replaced, in at least one way, by a reduced representation; such a representation can in particular be attained by successive decomposition.}

Indeed, let $\ideal M=[\ideal B_1^*\cdots\ideal B_l^*]$ be an arbitrary representation of $\ideal M$. We omit, in order, those $\ideal B_i^*$ which are absorbed into the least common multiple of the ideals retained. Since the remaining ideals still yield $\ideal M$, in the resulting representation
\[
 \ideal M=[\ideal B_1\cdots\ideal B_k]=[\ideal A_i,\ideal B_i]
\]
the first condition is satisfied, hence it is a shortest representation; and this condition remains satisfied if any $\ideal B_i$ is replaced by a proper divisor. The second condition, however, is always attainable by the theorem on the finite chain (Theorem I). For if
\[
 \ideal B_i\supset \ideal B_i'\supset \ideal B_i''\supset\cdots\supset \ideal B_i^{(\nu)}\supset\cdots,
 \qquad
 \ideal M=[\ideal A_i,\ideal B_i]=[\ideal A_i,\ideal B_i']=\cdots=[\ideal A_i,\ideal B_i^{(\nu)}]=\cdots,
\]
where every $\ideal B_i^{(\nu)}$ is a proper divisor of the one immediately preceding, then by that theorem the chain $\ideal B_i,\ideal B_i',\ldots,\ideal B_i^{(\nu)},\ldots$ must terminate after finitely many steps; hence in the representation $\ideal M=[\ideal A_i,\ideal B_i^{(\nu)}]$ the ideal $\ideal B_i^{(\nu)}$ can no longer be replaced by a proper divisor; and this holds a fortiori if $\ideal A_i$ is replaced by a proper divisor. Applying the procedure successively to each $\ideal B_i$, always forming the complement with the already reduced $\ideal B$, gives a reduced representation.\footnote{That such a representation is not uniquely defined by the given representation is shown by the previous example. For $(x^2,xy)=[(x),(x^2,xy,y^\lambda)]$, $\lambda\ge2$, besides $[(x),(x^2,y)]$ also $[(x),(x^2,ux+y)]$, for arbitrary $u$, is a reduced representation.}

In order to obtain such a representation successively, it remains to show that from the individual reduced representations
\[
 \ideal M=[\ideal B_1,\ideal C_1],\qquad
 \ideal C_1=[\ideal B_2,\ideal C_2],\quad\ldots,\quad
 \ideal C_{k-1}=[\ideal B_k,\ideal C_k]
\]
it follows that the resulting representation $\ideal M=[\ideal B_1\cdots\ideal B_k,\ideal C_k]$ is reduced. For this it suffices to show that from reduced representations
\[
 \ideal M=[\ideal B,\ideal C],\qquad \ideal C=[\ideal C_1,\ideal C_2]
\]
one obtains a reduced representation $\ideal M=[\ideal B,\ideal C_1,\ideal C_2]$. In fact, by hypothesis no $\ideal B$ lies in its complement; if this happened for a $\ideal C_i$, then, contrary to the hypothesis in the first representation, $\ideal C$ would be replaceable by a proper divisor, since by the second reduced representation $\ideal C_1$ and $\ideal C_2$ are proper divisors of $\ideal C$; so the representation is shortest. Moreover, by hypothesis no $\ideal B$ can be replaced by a proper divisor; if this were possible for a $\ideal C_i$, it would, contrary to the hypothesis, correspond to replacing $\ideal C$ by a proper divisor, since the representation for $\ideal C$ is reduced. Thus the lemma is proved.

\textbf{Definition II.} An ideal $\ideal M$ is called reducible if it can be represented as the least common multiple of two proper divisors; in the contrary case $\ideal M$ is called irreducible.

We now prove, by means of Theorem I on the finite chain and using reduced representations:

\textbf{Theorem II.} \emph{Every ideal is representable as the least common multiple of finitely many irreducible ideals.}\footnote{That such a representation is not unique in general is shown by the previous example: $(x^2,xy)=[(x),(x^2,ux+y)]$ for arbitrary $u$. The two components are irreducible for arbitrary $u$. Indeed, all divisors of $(x)$ are of the form $(x,g(y))$, where $g(y)$ denotes a polynomial in $y$; the least common multiple of any two therefore also has this form, hence is a proper divisor of $(x)$. The ideal $(x^2,ux+y)$ possesses only the single divisor $(x,y)$, and therefore is necessarily irreducible as well.}

For an arbitrary ideal $\ideal M$ is either irreducible; then $\ideal M=[\ideal M]$ is a representation of the kind required by Theorem II; or else $\ideal M=[\ideal B_1,\ideal C_1]$, where $\ideal B_1$ and $\ideal C_1$ are proper divisors of $\ideal M$, and the representation may, by Lemma I, be assumed reduced. For $\ideal C_1$ the same alternative holds: either it is irreducible, or there is a reduced representation $\ideal C_1=[\ideal B_2,\ideal C_2]$.

Continuing in this way, one obtains the series of reduced representations
\begin{equation}
 \ideal M=[\ideal B_1,\ideal C_1]=[\ideal B_1,\ideal B_2,\ideal C_2]=\cdots=[\ideal B_1,\ldots,\ideal B_n,\ideal C_n]=\cdots .
\tag{1}
\end{equation}
In the chain $\ideal C_1,\ideal C_2,\ldots,\ideal C_n,\ldots$ each $\ideal C_n$ is a proper divisor of the one immediately preceding, and therefore the chain terminates after finitely many steps; there is an index $n$ such that $\ideal C_n$ is irreducible. By Lemma I, moreover, the representation $\ideal M=[\ideal A_n,\ideal C_n]$ is reduced; $\ideal C_n$ therefore cannot lie in its complement $\ideal A_n$, and in the representation $\ideal M=[\ideal A_n,\ideal C_n]$ the ideal $\ideal A_n$ cannot be replaced by a proper divisor. Replacing, if necessary, $\ideal A_n$ by a proper divisor,\footnote{In fact the representation is also reduced with respect to $\ideal A_n$, as will be shown in \S~3 (Lemma IV) as the converse of Lemma I.} so that the representation becomes reduced, shows that every reducible ideal admits a reduced representation as the least common multiple of an irreducible ideal and a complementary ideal. In the series (1), therefore, all $\ideal B_i$ may be assumed irreducible without loss of generality; repeating the preceding argument gives the existence of an irreducible $\ideal C_n$, proving Theorem II.

\begin{center}
\S~3.\\[0.25em]
\textbf{Equality of the Number of Components in Two Different Decompositions into Irreducible Ideals.}
\end{center}

To prove equality of number, the reducibility or irreducibility of an ideal must first be expressed in terms of properties of its complement, as follows.

\textbf{Theorem III.}\footnote{Theorem III corresponds to the passage from modules to residue groups in the works of Schmeidler and Noether--Schmeidler (cf. the Introduction). The residue group corresponds to $\ideal C$, and the subgroups into which the residue group is decomposed correspond to $\ideal N_1$ and $\ideal N_2$.} \emph{Let the shortest representation $\ideal M=[\ideal A,\ideal C]$ be reduced with respect to $\ideal C$. Then the necessary and sufficient condition for $\ideal C$ to be reducible is the existence of two ideals $\ideal N_1$ and $\ideal N_2$ that are proper divisors of $\ideal M$, such that}
\srcnumdisplay[3.1em]{(2)}{\ideal N_1\equiv0(\ideal A);\qquad
 \ideal N_2\equiv0(\ideal A);\qquad
 [\ideal N_1,\ideal N_2]=\ideal M.}
\emph{It follows further: if conditions (2) are satisfied and $\ideal C$ is irreducible, then at least one $\ideal N_i$ is not a proper divisor of $\ideal M$; $\ideal N_i=\ideal M$.}

Let $\ideal C=[\ideal C_1,\ideal C_2]$, where $\ideal C_1,\ideal C_2$ are proper divisors of $\ideal C$. Then
\[
 \ideal M=[\ideal A,\ideal C]=[\ideal A,\ideal C_1,\ideal C_2]
       =[[\ideal A,\ideal C_1],[\ideal A,\ideal C_2]].
\]
Here the ideals $[\ideal A,\ideal C_i]$ are proper divisors of $\ideal M$, since otherwise $[\ideal A,\ideal C]$ would not be reduced with respect to $\ideal C$. Since divisibility by $\ideal A$ is also satisfied, condition (2) has been proved necessary. (The representation (2) is not reduced, since one $[\ideal A,\ideal C_i]$ can be replaced by $\ideal C_i$.)

Conversely, suppose (2) is satisfied. We form the ideals
\[
 \ideal C_1=(\ideal C,\ideal N_1),\qquad
 \ideal C_2=(\ideal C,\ideal N_2),\qquad
 \ideal C^*=[\ideal C_1,\ideal C_2].
\]
Then $\ideal C$ is divisible both by $\ideal C_1$ and by $\ideal C_2$, and hence also by their least common multiple $\ideal C^*$. To show that $\ideal C^*$ is divisible by $\ideal C$, let
\[
 f\equiv0(\ideal C^*);\quad f\equiv0(\ideal C_1);\quad f\equiv0(\ideal C_2),
 \quad\text{or equivalently}\quad f=c+n_1=\bar c+n_2,
\]
where $c,\bar c$ are elements of $\ideal C$ and $n_1,n_2$ are elements of $\ideal N_1,\ideal N_2$, respectively; in particular, $n_i$ is divisible by $\ideal N_i$. Thus the difference
\[
 g=c-\bar c=n_2-n_1
\]
is divisible both by $\ideal C$ and by $\ideal A$, hence by $\ideal M$. Since $n_1=n_2+m$, moreover, $n_1$ (and similarly $n_2$) is divisible both by $\ideal N_1$ and by $\ideal N_2$, hence by $\ideal M$. Therefore
\[
 f=c+m,\qquad f\equiv0(\ideal C),\qquad \ideal C^*=\ideal C.
\]
The ideals $\ideal C_1$ and $\ideal C_2$ are proper divisors of $\ideal C$; for if $\ideal C_1=(\ideal C,\ideal N_1)=\ideal C$, then $\ideal N_1$ would be divisible by $\ideal C$ and hence, because it is also divisible by $\ideal A$, would equal $\ideal M$, contrary to the hypothesis. Thus $\ideal C=[\ideal C_1,\ideal C_2]$ has been shown to be reducible; Theorem III is proved.

It may be noted that almost the same argument also proves the following.

\textbf{Lemma II.} \emph{If, in a shortest representation $\ideal M=[\ideal A,\ideal C]$, the ideal $\ideal C$ can be replaced by a proper divisor, then $\ideal C$ is reducible.}

Indeed, suppose
\[
 \ideal M=[\ideal A,\ideal C]=[\ideal A,\ideal C_1],
\]
and put
\[
 \ideal C^*=[\ideal C_1,(\ideal A,\ideal C)].
\]
Again $\ideal C$ is divisible by $\ideal C^*$. From $f\equiv0(\ideal C^*)$ it follows that
\[
 f=c_1=a+c.
\]
The difference $a=c_1-c$ is therefore divisible both by $\ideal A$ and by $\ideal C_1$, hence by $\ideal M$; thus $c_1=c+m$, $f\equiv0(\ideal C)$, and $\ideal C^*=\ideal C$. Since, by hypothesis, both $\ideal C_1$ and $(\ideal A,\ideal C)$ are proper divisors of $\ideal C$, the ideal $\ideal C=\ideal C^*$ has thereby been shown to be reducible.

An irreducible $\ideal C$ therefore cannot be replaced by a proper divisor.

Now let two different shortest representations of $\ideal M$ as the least common multiple of finitely many irreducible ideals be given:
\[
 \ideal M=[\ideal B_1\cdots\ideal B_k]=[\ideal D_1\cdots\ideal D_l].
\]
By the remark following Lemma II, these representations are at the same time reduced. We first prove:

\textbf{Lemma III.} \emph{For every complement $\ideal A_i=[\ideal B_1\cdots\ideal B_{i-1}\ideal B_{i+1}\cdots\ideal B_k]$ there is an ideal $\ideal D_j$ such that $\ideal M=[\ideal A_i,\ideal D_j]$.}

Indeed, put $\ideal M=[\ideal D_1,\ideal C_1]$, $\ideal C_1=[\ideal D_2,\ideal C_{12}]$, and so on. Then
\[
 \ideal M=[\ideal A_i,\ideal M]=[\ideal A_i,\ideal D_1,\ideal C_1]
 =[[\ideal A_i,\ideal D_1],[\ideal A_i,\ideal C_1]].
\]
Here, for $\ideal N_1=[\ideal A_i,\ideal D_1]$ and $\ideal N_2=[\ideal A_i,\ideal C_1]$, the conditions (2) of Theorem III are satisfied, since $\ideal M=[\ideal A_i,\ideal B_i]$ is reduced with respect to $\ideal B_i$ and the representation is a shortest one. Since $\ideal B_i$ was assumed irreducible, one $\ideal N_i$ must necessarily equal $\ideal M$.

If $\ideal N_1=\ideal M$, the lemma is proved. If $\ideal N_2=\ideal M$, then correspondingly
$\ideal M=[[\ideal A_i,\ideal D_2],[\ideal A_i,\ideal C_{12}]]$, where the same argument again shows that one component must equal $\ideal M$. Continuing in this way, either $\ideal M=[\ideal A_i,\ideal D_j]$ with $j<l$, or $\ideal M=[\ideal A_i,\ideal C_{1\ldots l-1}]$; since $\ideal C_{1\ldots l-1}=\ideal D_l$, the lemma is proved.

It now follows that:

\textbf{Theorem IV.} \emph{In two different shortest representations of an ideal as the least common multiple of irreducible ideals, the number of components is the same.}

Indeed, the lemma gives, for $i=1$,
\[
 \ideal M=[\ideal A_1,\ideal B_1]=[\ideal A_1,\ideal D_{j_1}]
 = [\ideal D_{j_1},\ideal B_2,\ldots,\ideal B_k].
\]
Now consider the two decompositions
\[
 \ideal M=[\ideal D_{j_1},\ideal B_2,\ldots,\ideal B_k]
          =[\ideal D_1,\ideal D_2,\ldots,\ideal D_l],
\]
and repeat the preceding argument with respect to the complement $\bar{\ideal A}_2=[\ideal D_{j_1},\ideal B_3,\ldots,\ideal B_k]$ of $\ideal B_2$. This gives
\[
 \ideal M=[\bar{\ideal A}_2,\ideal B_2]=[\bar{\ideal A}_2,\ideal D_{j_2}]
          =[\ideal D_{j_1},\ideal D_{j_2},\ideal B_3,\ldots,\ideal B_k],
\]
and, by continuing the procedure,
\[
 \ideal M=[\ideal D_{j_1},\ideal D_{j_2},\ldots,\ideal D_{j_k}].
\]
Since by hypothesis the representation $\ideal M=[\ideal D_1\cdots\ideal D_l]$ is a shortest one, so that no $\ideal D$ can be omitted, the distinct ideals among the $\ideal D_{j_i}$ must exhaust all the $\ideal D$; hence $k\ge l$. Interchanging the $\ideal B$ and $\ideal D$ throughout the lemma and the subsequent arguments gives correspondingly $l\ge k$, and therefore $k=l$, proving equality of number. It follows further that the ideals $\ideal D_{j_i}$ are all mutually distinct, since otherwise a shortest representation by the $\ideal D_{j_i}$ would have fewer than $k$ components; the notation may therefore be chosen so that $\ideal D_{j_i}=\ideal D_i$. For the same reason, all the intermediate representations $\ideal M=[\ideal D_{j_1}\cdots\ideal D_{j_r},\ideal B_{r+1},\ldots,\ideal B_k]$ are shortest and hence, by the remark following Lemma II, reduced.

Equality of number yields a converse of Lemma I:

\textbf{Lemma IV.} \emph{If, in a reduced representation, the components are combined into groups and their least common multiples are formed, then the resulting representation is reduced. In other words: from a reduced representation}
\[
 \ideal M=[\ideal C_{11}\cdots\ideal C_{1\mu_1};\;\ldots;\;
 \ideal C_{\sigma1}\cdots\ideal C_{\sigma\mu_\sigma}]
\]
\emph{it follows that $\ideal M=[\ideal N_1\cdots\ideal N_\sigma]=[\ideal N_i,\ideal L_i]$ is also reduced, if $\ideal N_i=[\ideal C_{i1}\cdots\ideal C_{i\mu_i}]$.}

First, $\ideal N_i$ cannot be absorbed into its complement $\ideal L_i$, since this is true of none of its divisors $\ideal C_{ij}$; the representation is therefore a shortest one. To show that $\ideal N_i$ cannot be replaced by a proper divisor, decompose the $\ideal C$ into their irreducible ideals $\ideal B$,\footnote{This always means a shortest, and therefore reduced, representation by the $\ideal B$.} so that the representations, reduced by Lemma I, become
\[
 \ideal M=[\ideal B_{11}\cdots\ideal B_{1\lambda_1};\;\ldots;\;
 \ideal B_{\sigma1}\cdots\ideal B_{\sigma\lambda_\sigma}];
 \qquad
 \ideal N_i=[\ideal B_{i1}\cdots\ideal B_{i\lambda_i}].
\]
Now let $\ideal M=[\ideal N_i^*,\ideal L_i]$ be reduced with respect to $\ideal N_i^*$, and let $\ideal N_i^*$ be a proper divisor of $\ideal N_i$. By Lemma II,
\[
 \ideal N_i=[\ideal N_i^*,(\ideal N_i,\ideal L_i)],
\]
and this representation is necessarily reduced with respect to $\ideal N_i^*$, since otherwise $\ideal N_i^*$ could also be replaced in $\ideal M$ by a proper divisor. If necessary, also replace $(\ideal N_i,\ideal L_i)$ by a proper divisor so that a reduced representation of $\ideal N_i$ results. If the two components of $\ideal N_i$ are now decomposed into irreducible ideals, the number $\lambda_i$ of irreducible ideals of $\ideal N_i$ is the sum of the numbers belonging to the components; the number of irreducible ideals corresponding to the proper divisor $\ideal N_i^*$ is therefore necessarily smaller than $\lambda_i$. But then the decomposition of $\ideal M=[\ideal N_i^*,\ideal L_i]$ into irreducible ideals would also yield fewer than $\sum_i\lambda_i$ ideals, contradicting equality of number. As the special case $\sigma=2$, it follows further that the representation $\ideal M=[\ideal N_i,\ideal L_i]$ is also reduced with respect to the complement $\ideal L_i$.

\begin{center}
\S~4.\\[0.25em]
\textbf{Primary Ideals. Uniqueness of the Associated Prime Ideals in Two Different Decompositions into Irreducible Ideals.}
\end{center}

What follows concerns the connection between primary and irreducible ideals.

\textbf{Definition III.} \emph{An ideal $\ideal Q$ is called primary if $a\cdot b\equiv0(\ideal Q)$ and $a\not\equiv0(\ideal Q)$ necessarily imply $b^\varkappa\equiv0(\ideal Q)$, where the exponent $\varkappa$ is finite.}

The definition may also be stated as follows: if a product $a\cdot b$ is divisible by $\ideal Q$, then either one factor is divisible by $\ideal Q$, or a power of each factor is. \emph{In particular, if $\varkappa$ is always equal to 1, the ideal is called a prime ideal.}

By virtue of the existence of bases, the definition of a primary (respectively prime) ideal yields the following formulation, which involves only products of ideals\footnote{As usual, the product $\ideal A\cdot\ideal B$ of two ideals means the ideal consisting of all elements $a\cdot b$ and their finite sums.}:

\textbf{Definition IIIa.} \emph{An ideal $\ideal Q$ is called primary if $\ideal A\cdot\ideal B\equiv0(\ideal Q)$ and $\ideal A\not\equiv0(\ideal Q)$ necessarily imply $\ideal B^\lambda\equiv0(\ideal Q)$. If $\lambda$ is always equal to 1, the ideal is called a prime ideal. Thus, for a prime ideal $\ideal P$, $\ideal A\cdot\ideal B\equiv0(\ideal P)$ and $\ideal A\not\equiv0(\ideal P)$ always imply $\ideal B\equiv0(\ideal P)$.}

Indeed, since Definition III is contained in IIIa as the special case $\ideal A=(a)$, $\ideal B=(b)$, every ideal primary according to IIIa is also primary according to III. Conversely, let $\ideal Q$ be primary according to III and suppose the hypothesis of IIIa is satisfied: $\ideal A\cdot\ideal B\equiv0(\ideal Q)$. Then either $\ideal A\equiv0(\ideal Q)$, or there is at least one element $a\equiv0(\ideal A)$ such that $a\cdot\ideal B\equiv0(\ideal Q)$ and $a\not\equiv0(\ideal Q)$. If $b_1,\ldots,b_r$ is an ideal basis of $\ideal B$, then Definition III, since $a\cdot b_i\equiv0(\ideal Q)$, gives
\[
 b_1^{\varkappa_1}\equiv0(\ideal Q);\quad\ldots;\quad
 b_r^{\varkappa_r}\equiv0(\ideal Q).
\]
Since
\[
 b=f_1b_1+\cdots+f_r b_r+n_1b_1+\cdots+n_r b_r,
\]
the product of any $\lambda=\varkappa_1+\cdots+\varkappa_r$ elements $b$ is therefore divisible by $\ideal Q$, proving that ideals primary according to III satisfy Definition IIIa. In the special case of prime ideals $\ideal P$, however, $a\cdot\ideal B\equiv0(\ideal P)$, and hence $a\cdot b\equiv0(\ideal P)$ for every $b\equiv0(\ideal B)$, together with $a\not\equiv0(\ideal P)$, imply $b\equiv0(\ideal P)$ and thus $\ideal B\equiv0(\ideal P)$. This proves the equivalence of the two definitions.

The connection between primary and prime ideals is established by observing that the totality $\ideal P$ of all elements $p$ having the property that some power of $p$ is divisible by $\ideal Q$ forms a prime ideal. First, it is clear that $\ideal P$ is an ideal, since, along with $p_1$ and $p_2$, the elements $ap_1$ and $p_1-p_2$ also have the stated property. By the same basis argument used in Definition IIIa, there is furthermore a number $\lambda$ such that $\ideal P^\lambda\equiv0(\ideal Q)$.

Now let
\[
 ab\equiv0(\ideal P),\qquad a\not\equiv0(\ideal P).
\]
Then, by the definition of $\ideal P$,
\[
 a^\lambda b^\lambda\equiv0(\ideal Q),\qquad a^\lambda\not\equiv0(\ideal Q),
\]
and therefore, by the definition of $\ideal Q$,
\[
 b^{\lambda\varkappa}\equiv0(\ideal Q),\qquad\text{and consequently}\qquad b\equiv0(\ideal P),
\]
which proves that $\ideal P$ is a prime ideal. The ideal $\ideal P$ is also characterized as the greatest common divisor of all ideals $\ideal B$ having the property that some power of $\ideal B$ is divisible by $\ideal Q$. For every such $\ideal B$ is divisible by $\ideal P$, and hence so is their greatest common divisor $\ideal D$. Conversely, $\ideal P$ itself is one of these ideals $\ideal B$, and hence is divisible by $\ideal D$, proving $\ideal P=\ideal D$. Thus $\ideal P$ is a prime ideal that is a divisor of $\ideal Q$ and some power of which is divisible by $\ideal Q$; this property determines it uniquely. For from
\[
 \ideal Q\equiv0(\ideal P),\qquad \ideal P^\lambda\equiv0(\ideal Q),
 \qquad
 \ideal Q\equiv0(\bar{\ideal P}),\qquad \bar{\ideal P}^{\mu}\equiv0(\ideal Q)
\]
it follows that
\[
 \ideal P^\lambda\equiv0(\bar{\ideal P}),\qquad
 \bar{\ideal P}^{\mu}\equiv0(\ideal P),
\]
and hence, by the property of prime ideals,
\[
 \ideal P\equiv0(\bar{\ideal P}),\qquad
 \bar{\ideal P}\equiv0(\ideal P),\qquad
 \ideal P=\bar{\ideal P}.
\]
In summary, we have:

\textbf{Theorem V.} \emph{For every primary ideal $\ideal Q$ there exists one and only one prime ideal $\ideal P$ that is a divisor of $\ideal Q$ and some power of which is divisible by $\ideal Q$; $\ideal P$ will be called the ``associated prime ideal''.}\footnote{That the converse does not hold is shown by the example $\ideal M=(x^2,xy)$. The prime ideal $(x)$ satisfies all the conditions, but $\ideal M$ is not primary.} \emph{$\ideal P$ is characterized as the greatest common divisor of all ideals $\ideal B$ having the property that some power of $\ideal B$ is divisible by $\ideal Q$. If $\varrho$ is the least number such that $\ideal P^\varrho\equiv0(\ideal Q)$, then $\varrho$ will be called the exponent of $\ideal Q$.}\footnote{That in general $\ideal P^\varrho=\ideal Q$ need not hold, as it does in the domain of rational integers and in that of algebraic integers, is shown by the example $\ideal Q=(x^2,y)$; $\ideal P=(x,y)$; $\ideal P^2=(x^2,xy,y^2)\equiv0(\ideal Q)$, but $\ideal Q\not\equiv0(\ideal P^2)$; thus $\ideal Q$ differs from $\ideal P^2$.}

We now prove, as the connection between primary and irreducible ideals:

\textbf{Theorem VI.} \emph{Every nonprimary ideal is reducible; in other words, every irreducible ideal is primary.}\footnote{That the converse does not hold here is shown, for example, by $\ideal Q=(x^2,xy,y^\lambda)=[(x^2,y),(x,y^\lambda)]$, where $\lambda\ge2$. Here $\ideal Q$ is primary but reducible. That $\ideal Q$ is primary follows from the fact that it contains all power products of degree $\lambda$ in $x,y$; thus, for every polynomial without a constant term, some power is divisible by $\ideal Q$. But if, in $a\cdot b\equiv0(\ideal Q)$, the polynomial $b$ contains a constant term---so that $b^\varkappa\not\equiv0(\ideal Q)$ for every $\varkappa$---then, since every homogeneous component of $a\cdot b$ is divisible by $\ideal Q$ because the basis polynomials of $\ideal Q$ are homogeneous, $a$ must be divisible by $\ideal Q$.}

Let $\ideal K$ be a nonprimary ideal, so that by Definition III there is at least one pair of elements $a,b$ such that
\srcnumdisplay[3.1em]{(3)}{a\cdot b\equiv0(\ideal K);\qquad
 a\not\equiv0(\ideal K);\qquad
 b^\varkappa\not\equiv0(\ideal K)\quad\text{for every }\varkappa.}
We now form the two ideals
\[
 \ideal L_0=(\ideal K,a);\qquad \ideal N_0=(\ideal K,b),
\]
which by (3) are proper divisors of $\ideal K$ and for which, again by (3),
\srcnumdisplay[3.1em]{(4)}{\ideal L_0\cdot\ideal N_0\equiv0(\ideal K).}
For the elements $f$ of the least common multiple $\ideal K_0=[\ideal L_0,\ideal N_0]$, the following alternative holds.

\emph{Either} from
\[
 f\equiv0(\ideal L_0);\qquad f\equiv0(\ideal N_0),
 \qquad\text{i.e.}\qquad f\equiv a_1\cdot b\;(\ideal K)
\]
there always follows a representation
\[
 f\equiv l_0\cdot b\;(\ideal K);\qquad l_0\equiv0(\ideal L_0);
\]
then by (4), $f\equiv0(\ideal K)$, hence $\ideal K_0\equiv0(\ideal K)$, and, since also $\ideal K\equiv0(\ideal K_0)$, $\ideal K=\ideal K_0$, which shows that $\ideal K$ is reducible.

\emph{Or} there is at least one $f\equiv0(\ideal K_0)$ for which no such $l_0$ exists. With the $a_1$ belonging to this $f$, we then form
\[
 \ideal L_1=(\ideal L_0,a_1)=(\ideal K,a,a_1);\qquad
 \ideal N_1=(\ideal K,b^2).
\]
Then, since $a_1\cdot b\equiv0(\ideal L_0)$, (4) also gives
\srcnumdisplay[3.1em]{(4')}{\ideal L_1\cdot\ideal N_1\equiv0(\ideal K);}
and $\ideal L_1$ is a proper divisor of $\ideal L_0$.

The same alternative holds for the elements $f$ of $\ideal K_1=[\ideal L_1,\ideal N_1]$.

\emph{Either} from
\[
 f\equiv0(\ideal L_1);\qquad f\equiv0(\ideal N_1);
\]
\[
 \text{i.e.}\qquad f\equiv a_2\cdot b^2\;(\ideal K)
 \quad\text{there always follows}\quad
 f\equiv l_1\cdot b^2;\qquad l_1\equiv0(\ideal L_1);
\]
and therefore by (4'), $\ideal K=\ideal K_1$.

\emph{Or} for at least one $f$ there is no such $l_1$, which leads to the formation of $\ideal L_2=(\ideal L_1,a_2)$ and $\ideal N_2=(\ideal K,b^{2^2})$, with $\ideal L_2\cdot\ideal N_2\equiv0(\ideal K)$, where $\ideal L_2$ is a proper divisor of $\ideal L_1$. Continuing in this way, we define in general
\[
 \ideal L_0=(\ideal K,a);\quad
 \ideal L_1=(\ideal L_0,a_1);\quad\ldots;\quad
 \ideal L_\nu=(\ideal L_{\nu-1},a_\nu);\quad\ldots,
\]
\[
 \ideal N_0=(\ideal K,b);\quad
 \ideal N_1=(\ideal K,b^2);\quad\ldots;\quad
 \ideal N_\nu=(\ideal K,b^{2^\nu});\quad\ldots,
\]
where the $a_i$ are defined by the existence of an $f$ such that
\[
 f\equiv0(\ideal L_{i-1});\qquad f\equiv0(\ideal N_{i-1}),
\]
\[
 \text{i.e.}\qquad f\equiv a_i\cdot b^{2^{i-1}}\;(\ideal K);
 \qquad\text{but}\qquad a_i\not\equiv0(\ideal L_{i-1}).
\]
It then follows in general that $\ideal L_i\cdot\ideal N_i\equiv0(\ideal K)$, that $\ideal N_i$ is a proper divisor of $\ideal K$ by (3), and that $\ideal L_i$ is a proper divisor of $\ideal L_{i-1}$. By Theorem I on the finite chain, the chain of the $\ideal L_i$ must therefore terminate after finitely many steps, say with $\ideal L_n$. Thus, for every $f\equiv0(\ideal L_n)$ and $f\equiv0(\ideal N_n)$, one has $f\equiv l_n\cdot b^{2^n}\;(\ideal K)$ with $l_n\equiv0(\ideal L_n)$, and consequently the preceding argument gives $\ideal K=[\ideal L_n,\ideal N_n]$, which shows that $\ideal K$ is reducible.

The \emph{uniqueness of the associated prime ideals} follows from what has just been proved, as follows.

Let
\[
 \ideal M=[\ideal B_1\cdots\ideal B_k]=[\ideal D_1\cdots\ideal D_k]
\]
be two shortest, hence reduced, representations of $\ideal M$ as the least common multiple of irreducible ideals, whose numbers of components agree by Theorem IV. By that theorem the intermediate representations occurring there (where, as noted there, the index $j_i=i$ may be put)
\[
 \ideal M=[\ideal D_1\cdots\ideal D_{i-1}\ideal B_i\ideal B_{i+1}\cdots\ideal B_k]
       =[\ideal D_1\cdots\ideal D_{i-1}\ideal D_i\ideal B_{i+1}\cdots\ideal B_k]
       =[\bar{\ideal U}_i,\ideal B_i]=[\bar{\ideal U}_i,\ideal D_i]
\]
are also shortest representations. Hence
\[
 \bar{\ideal U}_i\cdot\ideal B_i\equiv0(\ideal D_i),\qquad \bar{\ideal U}_i\not\equiv0(\ideal D_i),
 \qquad
 \bar{\ideal U}_i\cdot\ideal D_i\equiv0(\ideal B_i),\qquad \bar{\ideal U}_i\not\equiv0(\ideal B_i).
\]
Since by Theorem VI the irreducible ideals $\ideal B_i$ and $\ideal D_i$ are primary, it follows that there are two numbers $\lambda_i$ and $\mu_i$ such that
\srcnumdisplay[3.1em]{(5)}{\ideal B_i^{\lambda_i}\equiv0(\ideal D_i);\qquad
 \ideal D_i^{\mu_i}\equiv0(\ideal B_i).}
Let $\ideal P_i$ and $\bar{\ideal P}_i$ denote the associated prime ideals of $\ideal B_i$ and $\ideal D_i$, respectively; thus $\ideal P_i^{\varrho_i}\equiv0(\ideal B_i)$ and $\bar{\ideal P}_i^{\sigma_i}\equiv0(\ideal D_i)$. Then (5) gives
\[
 \ideal P_i^{\lambda_i\varrho_i}\equiv0(\bar{\ideal P}_i),
 \qquad
 \bar{\ideal P}_i^{\mu_i\sigma_i}\equiv0(\ideal P_i),
\]
and hence, by the property of prime ideals,
\[
 \ideal P_i\equiv0(\bar{\ideal P}_i),
 \qquad
 \bar{\ideal P}_i\equiv0(\ideal P_i),
 \qquad
 \ideal P_i=\bar{\ideal P}_i.
\]
This proves:

\textbf{Theorem VII.} \emph{In two different shortest representations of an ideal as the least common multiple of irreducible ideals, the associated prime ideals agree; equal prime ideals may occur among them,\footnote{This is shown, for example, by the example in footnote 19: $(x^2,xy,y^\lambda)=[(x^2,y),(x,y^\lambda)]$, where $\lambda\ge2$. The two associated prime ideals here are $(x,y)$.} and occur equally often in every decomposition. Consequently, the ideals themselves can be paired in at least one way so that, in each pair, a power of one ideal $\ideal B_i$ is divisible by the assigned ideal $\ideal D_i$, and conversely. Their numbers agree by Theorem IV.}\footnote{For the uniqueness of the ``isolated'' ideals among the irreducible ideals, cf. \S~7.}

\begin{center}
\S~5.\\[0.25em]
\textbf{Representation of an ideal as the least common multiple of greatest primary ideals. Uniqueness of the associated prime ideals.}
\end{center}

\textbf{Definition IV.} \emph{A shortest representation $\ideal M=[\ideal Q_1\cdots\ideal Q_\alpha]$ will be called a least common multiple of greatest primary ideals if all $\ideal Q$ are primary, but the least common multiple of two $\ideal Q$ is no longer primary.}

That at least one such representation always exists follows from the representation of $\ideal M$ as a least common multiple of irreducible ideals. These ideals are primary; either there is already a representation by greatest primary ideals, or the least common multiple of some two ideals is again primary. Since the number of ideals has thereby decreased by one, repetition of the procedure leads after finitely many steps to the desired representation.

This representation is reduced by Lemma IV. Conversely, every reduced representation by greatest primary ideals arises in this way, as the decomposition of the $\ideal Q$ into irreducible ideals shows.

In order to derive here from Theorem VII a corresponding uniqueness theorem, the connection with the associated prime ideals has to be investigated.

\textbf{Theorem VIII.} \emph{If the primary ideals $\ideal N_1,\ideal N_2,\ldots,\ideal N_\lambda$ all possess the same associated prime ideal $\ideal P$, then their least common multiple $\ideal Q=[\ideal N_1,\ideal N_2,\ldots,\ideal N_\lambda]$ is also primary and has $\ideal P$ as associated prime ideal. Conversely, if $\ideal Q=[\ideal N_1\cdots\ideal N_\lambda]$ is a reduced representation of the primary ideal $\ideal Q$, then all $\ideal N_i$ are primary and have, as associated prime ideal, the associated prime ideal $\ideal P$ of $\ideal Q$.}

For the proof of the first part, first observe that from $\ideal P^{\varrho_i}\equiv0(\ideal N_i)$ for every $i$ it also follows that $\ideal P^\tau\equiv0(\ideal Q)$, where $\tau$ denotes the largest of the indices $\varrho_i$. Since $\ideal P$ is also a divisor of $\ideal Q$, $\ideal P$ is necessarily the associated prime ideal if $\ideal Q$ is primary. From
\[
 \ideal A\ideal B\equiv0(\ideal Q),
 \qquad
 \ideal B^k\not\equiv0(\ideal Q)\quad\text{(for every }k)
\]
it follows that
\[
 \ideal B\not\equiv0(\ideal P),
 \quad\text{hence}\quad
 \ideal B^k\not\equiv0(\ideal N_i)\quad\text{(for every }k),
\]
and consequently $\ideal A\equiv0(\ideal N_i)$; hence $\ideal A\equiv0(\ideal Q)$. Thus $\ideal Q$ is proved primary, and $\ideal P$ the associated prime ideal.

Conversely, let first
\[
 \ideal Q=[\ideal N_1\cdots\ideal N_\lambda]=[\ideal N_i,\ideal L_i]
\]
be a shortest representation of $\ideal Q$ by primary ideals $\ideal N_i$, and let $\ideal P_i$ denote the corresponding associated prime ideals. From
\[
 \ideal L_i\cdot\ideal N_i\equiv0(\ideal Q),\qquad
 \ideal L_i\not\equiv0(\ideal Q)
\]
(because the representation is shortest) it follows that
\[
 \ideal N_i^{\sigma_i}\equiv0(\ideal Q),\qquad
 \text{or}\qquad
 \ideal P_i^{\varrho_i\sigma_i}\equiv0(\ideal Q).
\]
Since $\ideal P_i$ is at the same time a divisor of $\ideal Q$, $\ideal P_i$ therefore agrees, for every $i$, with the associated prime ideal $\ideal P$ of $\ideal Q$.

It remains to show that in every reduced representation $\ideal Q=[\ideal N_1\cdots\ideal N_\lambda]$ the $\ideal N_i$ are primary.\footnote{That the reduced representation is essential here is shown by the example of the non-reduced representation $\ideal Q=[x^2,xy,y^\lambda]=[(x^2,xy,y^2,yz),(x,y^\lambda)]$, where $\lambda\ge2$. Here $(x^2,xy,y^2,yz)=[(x^2,y),(x,y^2,z)]$ is not primary by what has just been proved; for the latter representation is a shortest one by primary ideals, but the associated prime ideals $(x,y)$ and $(x,y,z)$ are different. ($\ideal Q$ is primary by footnote 19.)} To this end, decompose the $\ideal N_i$ into their irreducible ideals $\ideal B$; in the resulting shortest representation $\ideal Q=[\ideal B_1\cdots\ideal B_s]$, every ideal is primary and, by what has just been proved, has $\ideal P$ as associated prime ideal. Then by the first part of the theorem proved above the same holds for each $\ideal N_i$, completing the proof of Theorem VIII.

\textbf{Addendum.} It also follows that a prime ideal is necessarily irreducible. For from the reduced representation $\ideal P=[\ideal N,\ideal K]$ one obtains by Theorem VIII $\ideal N\equiv0(\ideal P)$; hence, since $\ideal P\equiv0(\ideal N)$, also $\ideal P=\ideal N$, and similarly $\ideal P=\ideal K$. The irreducibility of $\ideal P$ also follows directly: for from $\ideal P=[\ideal N,\ideal K]$ follows $\ideal N\ideal K\equiv0(\ideal P)$, $\ideal N\not\equiv0(\ideal P)$, $\ideal K\not\equiv0(\ideal P)$, contrary to the defining property of a prime ideal.

Now let
\[
 \ideal M=[\ideal Q_1\cdots\ideal Q_\alpha]=[\bar{\ideal Q}_1\cdots\bar{\ideal Q}_\beta]
\]
be two reduced representations of $\ideal M$ as the least common multiple of greatest primary ideals. Decomposing the $\ideal Q$ into irreducible ideals $\ideal B$, and the $\bar{\ideal Q}$ correspondingly into irreducible ideals $\ideal D$, gives two reduced representations of $\ideal M$ as the least common multiple of irreducible ideals; by Theorem VII, both the number of components and the associated prime ideals coincide. By Theorem VIII, all irreducible ideals $\ideal B$ occurring for a fixed $\ideal Q_i$ have the same associated prime ideal $\ideal P_i$; whereas the $\ideal P_j$ associated to $\ideal Q_j$ is necessarily different from this, since otherwise, by Theorem VIII, there would be no representation by greatest primary ideals. Hence the number $\alpha$ of the $\ideal Q$ is equal to the number of different associated prime ideals $\ideal P$ of the $\ideal B$; these different $\ideal P$ form the associated prime ideals of the $\ideal Q$. The same holds for the $\bar{\ideal Q}$ with respect to their decomposition into the $\ideal D$. Theorem VII therefore gives equality of number for the $\ideal Q$ and $\bar{\ideal Q}$, and agreement of their associated prime ideals. At the same time it is seen that the grouping of the irreducible ideals into greatest primary ideals described at the beginning of the paragraph consists in combining all, and only those, with the same associated prime ideal. Theorem VIII further shows the irreducibility property of the greatest primary ideals: they admit no reduced representation as the least common multiple of greatest primary ideals.

In summary:

\textbf{Theorem IX.} \emph{In two reduced representations of an ideal as the least common multiple of greatest primary ideals, the number of components and the associated prime ideals, all of which are distinct from each other, coincide. In other words, to each $\ideal Q$ there can be assigned one and only one $\bar{\ideal Q}$ such that a power of $\ideal Q$ is divisible by $\bar{\ideal Q}$, and conversely.}\footnote{An example of different representations is the one given in footnote 12 for Theorem II: $(x^2,xy)=[(x),(x^2,\mu x+y)]$ for arbitrary $\mu$. Since the associated prime ideals $\ideal P_1=(x)$ and $\ideal P_2=(x,y)$ are distinct, these are greatest primary ideals. -- For the uniqueness of the ``isolated'' greatest primary ideals, compare \S~7.} \emph{The $\ideal Q$ and $\bar{\ideal Q}$ have the irreducibility property with respect to decomposition into greatest primary ideals.}

\textbf{Addendum.} It should be noted that Theorem IX remains essentially valid if one assumes only a shortest representation instead of a reduced one. If, say,
\[
 \ideal M=[\ideal Q_1\cdots\ideal Q_i^*\cdots\ideal Q_\alpha]
\]
is reduced with respect to $\ideal Q_i^*$, and $\ideal Q_i^*$ is a proper divisor of $\ideal Q_i$, while $\ideal L_i$ is the complement of $\ideal Q_i$, then by Lemma IV
\[
 \ideal Q_i=[\ideal Q_i^*;(\ideal L_i,\ideal Q_i)],
\]
and this representation is reduced with respect to $\ideal Q_i^*$. By Theorem VIII -- replacing $(\ideal L_i,\ideal Q_i)$ by a proper divisor if necessary when applying the theorem -- $\ideal Q_i^*$ is therefore primary and has the same associated prime ideal $\ideal P_i$ as $\ideal Q_i$. Continuing the procedure shows that to every such representation there can be assigned a reduced representation by greatest primary ideals in such a way that the number of components and the associated prime ideals coincide. Thus even for shortest representations it holds that in two different representations the number of components and the associated prime ideals coincide.

The associated, mutually distinct prime ideals thereby uniquely defined will be called simply ``the associated prime ideals of $\ideal M$.''




\begin{center}
\S~6.\\[0.25em]
\textbf{Unique representation of an ideal as the least common multiple of relatively-prime irreducible ideals.}
\end{center}

\textbf{Definition V.} \emph{An ideal $\ideal R$ is called relatively prime to $\ideal S$ if $\ideal T\cdot\ideal R\equiv0(\ideal S)$ necessarily implies $\ideal T\equiv0(\ideal S)$. If both $\ideal R$ is relatively prime to $\ideal S$ and $\ideal S$ is relatively prime to $\ideal R$, then $\ideal R$ and $\ideal S$ are called mutually relatively prime.}\footnote{The condition of being relatively prime is not symmetric. For example, $\ideal R=(x^2,y)$ is relatively prime to $\ideal S=(x)$; but $\ideal S$ is not relatively prime to $\ideal R$, since $\ideal S^2\equiv0(\ideal R)$, whereas $\ideal S\not\equiv0(\ideal R)$.}

\emph{An ideal is called relatively-prime irreducible if it cannot be represented as the least common multiple of mutually relatively prime proper divisors.}

In particular, if instead of $\ideal T$ one takes the greatest common divisor $\ideal T_0$ of all $\ideal T$ for which $\ideal T\ideal R\equiv0(\ideal S)$, then also $\ideal T_0\ideal R\equiv0(\ideal S)$ and $\ideal S\equiv0(\ideal T_0)$. Thus $\ideal T_0=\ideal S$ when $\ideal R$ is relatively prime to $\ideal S$; and $\ideal T_0$ is a proper divisor of $\ideal S$ when $\ideal R$ is not relatively prime to $\ideal S$.\footnote{The $\ideal T_0$ so defined agrees with Lasker's ``residual module''; and, in its extension to number modules in place of ideals, with Dedekind's ``quotient'' of two modules. Lasker, loc. cit., p. 49; Dedekind (Zahlentheorie), p. 504.}

The uniqueness proof rests on:

\textbf{Theorem X.} \emph{1. If $\ideal R$ is relatively prime to the ideals $\ideal S_1,
\ldots,\ideal S_\lambda$, then $\ideal R$ is also relatively prime to their least common multiple $\ideal S$.}

\emph{2. If the ideals $\ideal S_1,
\ldots,\ideal S_\lambda$ are relatively prime to $\ideal R$, then their least common multiple $\ideal S$ is also relatively prime to $\ideal R$.}

\emph{3. If $\ideal R$ is relatively prime to $\ideal S$ and $\ideal S=[\ideal S_1\cdots\ideal S_\lambda]$ is a reduced representation for $\ideal S$, then $\ideal R$ is also relatively prime to each $\ideal S_i$.}

\emph{4. If $\ideal S$ is relatively prime to $\ideal R$, then every divisor $\ideal S_i$ of $\ideal S$ is also relatively prime to $\ideal R$.}

1. Since $\ideal T\ideal R\equiv0(\ideal S)$ necessarily gives $\ideal T\ideal R\equiv0(\ideal S_i)$, the hypothesis gives $\ideal T\equiv0(\ideal S_i)$, and hence $\ideal T\equiv0(\ideal S)$.

2. Put
\[
 \ideal C_1=[\ideal S_2\cdots\ideal S_\lambda],\quad
 \ideal C_{12}=[\ideal S_3\cdots\ideal S_\lambda],\quad \ldots,\quad
 \ideal C_{12\ldots\lambda-1}=\ideal S_\lambda.
\]
From $\ideal T\cdot\ideal S\equiv0(\ideal R)$ it follows that $\ideal T\cdot\ideal S_1\cdot\ideal C_1\equiv0(\ideal R)$; hence, by the hypothesis, $\ideal T\cdot\ideal C_1\equiv0(\ideal R)$. It follows in turn that $\ideal T\cdot\ideal S_2\cdot\ideal C_{12}\equiv0(\ideal R)$, hence $\ideal T\cdot\ideal C_{12}\equiv0(\ideal R)$, and finally $\ideal T\cdot\ideal C_{12\ldots\lambda-1}=\ideal T\cdot\ideal S_\lambda\equiv0(\ideal R)$; hence $\ideal T\equiv0(\ideal R)$.

3. Let $\ideal C_i$ denote the complement of $\ideal S_i$. From $\ideal T_0\cdot\ideal R\equiv0(\ideal S_i)$, together with $\ideal C_i\cdot\ideal R\equiv0(\ideal S_i)$, it follows that $[\ideal T_0,\ideal C_i]\cdot\ideal R\equiv0(\ideal S)$. If $\ideal T_0$ were a proper divisor of $\ideal S_i$, then, since $\ideal S=[\ideal S_i,\ideal C_i]$ is reduced, $[\ideal T_0,\ideal C_i]$ would be a proper divisor of $\ideal S$, contrary to the hypothesis.

4. From $\ideal T\ideal S_i\equiv0(\ideal R)$, where $\ideal T$ is a proper divisor of $\ideal R$, it follows also that $\ideal T\ideal S\equiv0(\ideal R)$; hence $\ideal T$ would be a proper divisor of $\ideal R$, contrary to the hypothesis.

Definition V shows in particular that every ideal is relatively prime to the unit ideal $\ideal D$ consisting of all elements of $\Sigma$;\footnote{$\ideal D$ plays the role of the unit only with respect to divisibility and least common multiple, not with respect to product formation. For instance, $\ideal D=(x)$ for the domain of all integral polynomials in $x$ without constant term; and $\ideal D=(2)$ for the domain of all even numbers.} the following theorems, however, apply only to ideals distinct from $\ideal D$.

From Theorem X one first obtains:

\textbf{Lemma V.} \emph{Every representation of an ideal as the least common multiple of mutually relatively prime ideals distinct from $\ideal D$ is reduced.}

Let
\[
 \ideal M=[\ideal R_1\ideal R_2\cdots\ideal R_\sigma]=[\ideal R_i,
\ideal L_i]
\]
be such a representation. By Theorem X, 2, $\ideal L_i$ is also relatively prime to $\ideal R_i$; hence $\ideal R_i$ cannot be contained in $\ideal L_i$, and the representation is shortest. For from $\ideal L_i\equiv0(\ideal R_i)$, with $\ideal L_i$ relatively prime to $\ideal R_i$, it would follow from $\ideal D\ideal L_i\equiv0(\ideal R_i)$ that $\ideal D\equiv0(\ideal R_i)$, hence $\ideal R_i=\ideal D$, contrary to the hypothesis.

If now $\ideal R_i$ could be replaced by a proper divisor $\ideal R_i^*$, then by Lemma II $\ideal R_i$ would be reducible and
\[
 \ideal R_i=[\ideal R_i^*,(\ideal R_i,
\ideal L_i)].
\]
If necessary, replace $(\ideal R_i,
\ideal L_i)$ by a proper divisor $(\ideal R_i,
\ideal L_i)^*$ and likewise replace $\ideal R_i^*$ by a proper divisor so that a reduced representation of $\ideal R_i$ is obtained. Then Theorem X, 3 shows that $\ideal L_i$ is also relatively prime to $(\ideal R_i,
\ideal L_i)^*$, and this ideal is distinct from $\ideal D$ because the original representation was shortest. But $(\ideal R_i,
\ideal L_i)^*$ is contained in $(\ideal R_i,
\ideal L_i)$ and $(\ideal R_i,
\ideal L_i)$ is contained in $\ideal L_i$, giving the contradiction.

Theorem X further gives the following theorem, which mediates the connection with the associated prime ideals.

\textbf{Theorem XI.} \emph{If $\ideal R$ is relatively prime to $\ideal S$ and $\ideal S$ is distinct from $\ideal D$, then no associated prime ideal of $\ideal R$\footnote{The associated prime ideals of $\ideal R$ were defined at the end of \S~5.} is divisible by an associated prime ideal of $\ideal S$. Conversely, if no such divisibility holds, then $\ideal R$ is relatively prime to $\ideal S$, and of course $\ideal S$ is distinct from $\ideal D$.}

Let
\[
 \ideal R=[\ideal Q_1\cdots\ideal Q_\alpha],
 \qquad
 \ideal S=[\ideal Q_1^*\cdots\ideal Q_\beta^*]
\]
be reduced representations of $\ideal R$ and $\ideal S$ by greatest primary ideals, and let $\ideal P_1,
\ldots,\ideal P_\alpha$ and $\ideal P_1^*,
\ldots,\ideal P_\beta^*$ be their associated prime ideals. We prove the assertion in the form: if some $\ideal P$ is divisible by some $\ideal P^*$, then $\ideal R$ cannot be relatively prime to $\ideal S$, and conversely.

Suppose then that $\ideal P_\mu\equiv0(\ideal P_\nu^*)$, and consequently also $\ideal Q_\mu\equiv0(\ideal P_\nu^*)$. By the definition of $\ideal P_\nu^*$ this gives $\ideal Q_\mu^{\sigma_\nu}\equiv0(\ideal Q_\nu^*)$, hence also $\ideal R^{\sigma_\nu}\equiv0(\ideal Q_\nu^*)$. Let $\ideal R^\tau$ be the lowest power of $\ideal R$ divisible by $\ideal Q_\nu^*$. If $\tau=1$, then $\ideal R$ is divisible by $\ideal Q_\nu^*$; but since $\ideal S\ne\ideal D$, $\ideal R$ is not relatively prime to $\ideal Q_\nu^*$. If $\tau\ge2$, then $\ideal R^{\tau-1}\ideal R\equiv0(\ideal Q_\nu^*)$ while $\ideal R^{\tau-1}\not\equiv0(\ideal Q_\nu^*)$, so again $\ideal R$ is not relatively prime to $\ideal Q_\nu^*$;\footnote{$\ideal R^0$ is not defined, because $\Sigma$ need not contain a unit; hence the case $\tau=1$ had to be dealt with first. The case $\tau=0$ is also excluded, if $\Sigma$ has a unit, by the hypothesis on $\ideal S$.} in both cases Theorem X, 3 then implies that $\ideal R$ is not relatively prime to $\ideal S$.

Conversely, if $\ideal R$ is not relatively prime to $\ideal S$, then by Theorem X, 1 it is not relatively prime to at least one $\ideal Q_\nu^*$. Thus
\[
 \ideal T_0\ideal R\equiv0(\ideal Q_\nu^*),
 \qquad
 \ideal T_0\not\equiv0(\ideal Q_\nu^*),
\]
and since $\ideal Q_\nu^*$ is primary,
\[
 \ideal R^\tau\equiv0(\ideal Q_\nu^*),\qquad
 \text{hence also}\quad
 \ideal Q_1^\tau\cdots\ideal Q_\alpha^\tau\equiv0(\ideal Q_\nu^*).
\]
For the associated prime ideals this implies
\[
 \ideal P_1^{\tau\varrho_1}\cdots\ideal P_\alpha^{\tau\varrho_\alpha}\equiv0(\ideal P_\nu^*),
\]
and therefore, by the property of prime ideals, $\ideal P_\nu^*$ is contained in at least one of the $\ideal P$. This proves Theorem XI.

Theorems X and XI now give the existence\footnote{The existence of the decomposition can also be proved directly, in exact analogy with the proof in \S~2 of the existence of a decomposition into finitely many irreducible ideals.} and uniqueness of decomposition into relatively-prime irreducible ideals as follows.

Let $\ideal M=[\ideal Q_1\cdots\ideal Q_\alpha]$ be a reduced (or at least shortest) representation of $\ideal M$ by greatest primary ideals, and let $\ideal P_1,
\ldots,\ideal P_\alpha$ be the associated prime ideals. We collect the $\ideal P$ into groups so that no ideal of one group is divisible by an ideal of a different group, while a single group cannot be split into two subgroups both having this property.

To construct such a grouping, note that by definition an individual group $G$ must contain, together with each ideal $\ideal P$, all its divisors and multiples occurring among $\ideal P_1,\ldots,\ideal P_\alpha$ (that is, all ideals divisible by $\ideal P$). Let $\ideal P^{(i_1)}$ denote all multiples of $\ideal P$, $\ideal P_{j_1}^{(i_1)}$ all divisors of $\ideal P^{(i_1)}$, $\ideal P_{j_1}^{(i_1i_2)}$ all multiples of $\ideal P_{j_1}^{(i_1)}$, and so on; in general, let $\ideal P_{j_1\ldots j_{\lambda-1}}^{(i_1\ldots i_\lambda)}$ denote all multiples of $\ideal P_{j_1\ldots j_{\lambda-1}}^{(i_1\ldots i_{\lambda-1})}$, and $\ideal P_{j_1\ldots j_\lambda}^{(i_1\ldots i_\lambda)}$ all divisors of $\ideal P_{j_1\ldots j_{\lambda-1}}^{(i_1\ldots i_\lambda)}$. Since only finitely many ideals $\ideal P$ are involved altogether, the procedure must terminate after finitely many steps, that is, must cease to produce an ideal different from all preceding ones. The totality of ideals $\ideal P$ so obtained does indeed form a group $G$ with the required properties.

For by definition $G$ contains, together with each $\ideal P_{j_1\ldots j_{\lambda-1}}^{(i_1\ldots i_{\lambda-1})}$, all multiples $\ideal P_{j_1\ldots j_{\lambda-1}}^{(i_1\ldots i_\lambda)}$, and, together with each $\ideal P_{j_1\ldots j_{\lambda-1}}^{(i_1\ldots i_\lambda)}$, all divisors $\ideal P_{j_1\ldots j_\lambda}^{(i_1\ldots i_\lambda)}$. These include all divisors of $\ideal P_{j_1\ldots j_{\lambda-1}}^{(i_1\ldots i_{\lambda-1})}$, while the multiples of $\ideal P_{j_1\ldots j_{\lambda-1}}^{(i_1\ldots i_\lambda)}$ are themselves again among the $\ideal P_{j_1\ldots j_{\lambda-1}}^{(i_1\ldots i_\lambda)}$. Thus the ideals not contained in $G$ can be neither divisors nor multiples of those contained in $G$.

The group $G$ also satisfies the irreducibility condition. For if a division into two subgroups $G^{(1)}$ and $G^{(2)}$ were possible, and $G^{(1)}$ contained, say, $\ideal P_{j_1\ldots j_\lambda}^{(i_1\ldots i_\lambda)}$ (respectively $\ideal P_{j_1\ldots j_{\lambda-1}}^{(i_1\ldots i_\lambda)}$), it would also contain all preceding ideals, since these are alternately multiples and divisors (respectively divisors and multiples), hence also $\ideal P$ and therefore the entire group $G$. Applying the same procedure to the ideals not contained in $G$ yields a grouping $G_1,\ldots,G_\sigma$ of all the $\ideal P$ with the required properties. Such a grouping is unique: for if $G'_1,\ldots,G'_\tau$ were a second grouping and $\ideal P_{j_1\ldots j_\lambda}^{(i_1\ldots i_\lambda)}$ (respectively $\ideal P_{j_1\ldots j_{\lambda-1}}^{(i_1\ldots i_\lambda)}$) were an element of $G'_i$, then by the preceding argument $G'_i$ would contain the entire group $G$ and hence, by the irreducibility property, would be identical with $G$.

Let $\ideal Q_{i\mu}$, where $\mu$ runs from $1$ to $\lambda_i$, denote the ideals $\ideal Q$ whose associated prime ideals are collected in a group $G_i$. Since the $\ideal P$ are all distinct, the primary ideals $\ideal Q_{i\mu}$ belonging to a shortest representation are uniquely determined. Put
\[
 \ideal R_i=[\ideal Q_{i1}\cdots\ideal Q_{i\lambda_i}],\qquad
 \text{then}\quad
 \ideal M=[\ideal R_1\cdots\ideal R_\sigma].
\]
Theorem XI remains applicable even if $\ideal R_i$ is only a shortest representation, since by the addendum to Theorem IX its associated prime ideals are already well-defined. Because no associated prime ideal of $\ideal R_i$ is divisible by one of $\ideal R_j$, and conversely, Theorem XI shows that $\ideal R_i$ and $\ideal R_j$ are mutually relatively prime; and each $\ideal R_i$ is relatively-prime irreducible by the irreducibility property of the group $G_i$. By Lemma V the representation is reduced. Conversely, every decomposition into relatively-prime irreducible ideals leads, by Theorem XI, to the grouping just described.

Now let
\[
 \ideal M=[\bar{\ideal R}_1\cdots\bar{\ideal R}_\tau]
\]
be a second representation of $\ideal M$ by relatively-prime irreducible ideals. It is reduced by Lemma V. Decomposing the $\ideal R$ into greatest primary ideals shows that the associated prime ideals agree, and hence the groupings of these prime ideals agree. Thus $\tau=\sigma$, and the notation may be chosen so that $\ideal R_i$ and $\bar{\ideal R}_i$ belong to the same group. If
\[
 \ideal M=[\ideal R_i,
\ideal L_i]=[\bar{\ideal R}_i,
\bar{\ideal L}_i]
\]
is the representation by ideal and complement, then, since $\bar{\ideal R}_i$ is assigned to the same group as $\ideal R_i$, Theorem XI gives that $\ideal L_i$ is relatively prime to $\bar{\ideal R}_i$, and $\bar{\ideal L}_i$ is relatively prime to $\ideal R_i$. From
\[
 \ideal R_i\ideal L_i\equiv0(\bar{\ideal R}_i),
 \qquad
 \bar{\ideal R}_i\bar{\ideal L}_i\equiv0(\ideal R_i)
\]
therefore follows
\[
 \ideal R_i\equiv0(\bar{\ideal R}_i),
 \qquad
 \bar{\ideal R}_i\equiv0(\ideal R_i),
 \qquad
 \ideal R_i=\bar{\ideal R}_i.
\]
Thus:

\textbf{Theorem XII.} \emph{Every ideal can be represented uniquely as the least common multiple of finitely many mutually relatively prime and relatively-prime irreducible ideals.}

\begin{center}
\S~7.\\[0.25em]
\textbf{Uniqueness of the Isolated Ideals.}
\end{center}

\textbf{Definition VI.} \emph{If the shortest representation $\ideal M=[\ideal R,
\ideal L]$ is reduced with respect to $\ideal L$, then $\ideal R$ is called an isolated ideal if no associated prime ideal of $\ideal R$ is contained in an associated prime ideal of $\ideal L$; equivalently, if $\ideal L$ is relatively prime to $\ideal R$.}

Then the representation $\ideal M=[\ideal R,
\ideal L]$ satisfies the conditions of the representation $\ideal M=[\ideal R_i,
\ideal L_i]$ in Lemma V, and Lemma V proves:

\textbf{Lemma VI.} \emph{If $\ideal R$ is an isolated ideal of the shortest representation $\ideal M=[\ideal R,
\ideal L]$, reduced with respect to $\ideal L$, then the representation is also reduced with respect to $\ideal R$.}

Thus isolated ideals occur only in reduced representations. The associated prime ideals arising in the decompositions of $\ideal R$ and $\ideal L$ into irreducible ideals complement one another to the uniquely determined associated prime ideals arising in the corresponding decomposition of $\ideal M$. Hence no associated prime ideal belonging to the decomposition of $\ideal R$ into irreducible ideals is contained in the remaining associated prime ideals of $\ideal M$. Conversely, if this condition is satisfied and $\ideal R$ occurs in at least one representation $\ideal M=[\ideal R,
\ideal L]$ reduced with respect to $\ideal R$, and hence also in a reduced representation $\ideal M=[\ideal R,
\ideal L^*]$, then $\ideal R$ is isolated by Definition VI. This gives the following definition, independent of the particular complement $\ideal L$.

\textbf{Definition VIa.} \emph{$\ideal R$ is called an isolated ideal if the prime ideals belonging to its decomposition into irreducible ideals are not contained in the other associated prime ideals belonging to the corresponding decomposition of $\ideal M$, and if $\ideal R$ occurs in at least one representation $\ideal M=[\ideal R,
\ideal L]$ reduced with respect to $\ideal R$.}\footnote{If one introduced ordinary associated prime ideals (end of \S~5), then one would have to add the special condition that those of $\ideal L$ are all distinct from those of $\ideal R$. With Definition VIa, the representation no longer has to be assumed reduced with respect to the complement.}

Suppose
\[
 \ideal M=[\ideal R,
\ideal L]=[\bar{\ideal R},
\bar{\ideal L}]
\]
are two representations of $\ideal M$ by isolated ideals $\ideal R$ and $\bar{\ideal R}$, with complements $\ideal L$ and $\bar{\ideal L}$, such that the associated prime ideals of $\ideal R$ and $\bar{\ideal R}$ agree. Replace $\ideal L$ and $\bar{\ideal L}$ by divisors $\ideal L^*$ and $\bar{\ideal L}^{*}$ so that the representations become reduced. Then the associated prime ideals of $\ideal L^*$ and $\bar{\ideal L}^{*}$ also agree; by Theorem XI, $\ideal L^*$ is relatively prime to $\bar{\ideal R}$ and $\bar{\ideal L}^{*}$ is relatively prime to $\ideal R$. Since
\[
 \ideal R\ideal L^*\equiv0(\bar{\ideal R}),
 \qquad
 \bar{\ideal R}\bar{\ideal L}^{*}\equiv0(\ideal R),
\]
we obtain
\[
 \ideal R\equiv0(\bar{\ideal R}),
 \qquad
 \bar{\ideal R}\equiv0(\ideal R),
 \qquad
 \ideal R=\bar{\ideal R}.
\]
Thus isolated ideals are uniquely determined by their associated prime ideals. This sharpens Theorems VII and IX on decompositions into irreducible, respectively greatest primary, ideals; by the addendum to Theorem IX it is enough to assume a shortest representation. In summary:

\textbf{Theorem XIII.} \emph{In every shortest representation of an ideal as the least common multiple of irreducible, respectively greatest primary, ideals, the isolated irreducible, respectively greatest primary, ideals are uniquely determined; the ambiguity concerns only the non-isolated irreducible, respectively greatest primary, ideals.}\footnote{For ideals in polynomial domains, this theorem was already stated without proof by Macaulay in the case of decomposition into greatest primary ideals; his definition of isolated and non-isolated (imbedded) primary ideals can be regarded as an irrational form of the one given below.} \emph{In general, isolated ideals are uniquely determined by their associated prime ideals.}

If in such a shortest representation by irreducible, respectively greatest primary, ideals the ideals $\ideal B_i$, respectively $\ideal Q_j$, are non-isolated, then by definition their complements $\ideal U_i$, respectively $\ideal L_j$, are divisible by $\ideal P_i$, respectively $\ideal P_j$. Hence
\[
 \ideal U_i^{\varrho_i}\equiv0(\ideal B_i)
 \qquad \text{respectively} \qquad
 \ideal L_j^{\sigma_j}\equiv0(\ideal Q_j).
\]
Conversely, if these relations hold, then $\ideal P_i$, respectively $\ideal P_j$, is contained in at least one associated prime ideal of the complement, and hence, by the addendum to Theorem IX, in an associated prime ideal of the divisor $\ideal L_j^*$ of $\ideal L_j$ that supplies a representation reduced with respect to $\ideal L_j^*$. Thus $\ideal B_i$, respectively $\ideal Q_j$, is non-isolated. In particular, irreducible ideals $\ideal B_i$ whose associated prime ideal occurs more than once in the decomposition of $\ideal M$ are always non-isolated. \emph{Non-isolated primary ideals are therefore also characterized by the fact that a power of every complement is divisible by them; isolated primary ideals by the fact that this cannot occur.}

\begin{center}
\S~8.\\[0.25em]
\textbf{Unique Representation of an Ideal as a Product of Coprime-Irreducible Ideals.}
\end{center}

If the underlying ring domain $\Sigma$ has a unit, that is, an element $\varepsilon$ such that $\varepsilon a=a$ for every element of $\Sigma$,\footnote{Because multiplication is commutative, $\Sigma$ can of course have at most one unit: for any two units $\varepsilon_1$ and $\varepsilon_2$, one has $\varepsilon_1\varepsilon_2=\varepsilon_2=\varepsilon_1$.} then coprime ideals can be defined as follows.

\textbf{Definition VIII.} \emph{Two ideals $\ideal R$ and $\ideal S$ are called coprime if their greatest common divisor is the unit ideal $\ideal D=(\varepsilon)$ consisting of all elements of $\Sigma$. An ideal is called coprime-irreducible if it cannot be represented as the least common multiple of pairwise coprime ideals.}

Note that two coprime ideals are always mutually relatively prime. By definition there are elements $r\equiv0(\ideal R)$ and $s\equiv0(\ideal S)$ with $\varepsilon=r+s$. From $\ideal T\ideal R\equiv0(\ideal S)$ it follows that $\ideal T r\equiv0(\ideal S)$, hence $\ideal T\varepsilon=\ideal T\equiv0(\ideal S)$; correspondingly, from $\bar{\ideal T}\ideal S\equiv0(\ideal R)$ it follows that $\bar{\ideal T}\equiv0(\ideal R)$.\footnote{The converse is false: for example, $\ideal R=(x)$ and $\ideal S=(y)$ are mutually relatively prime, but not coprime.} Thus Lemma V implies that every representation by pairwise coprime ideals is also reduced.

The uniqueness proof rests on the analogue of Theorem X:

\textbf{Theorem XIV.} \emph{If $\ideal R$ is coprime to each of the ideals $\ideal S_1,
\ldots,\ideal S_\lambda$, then $\ideal R$ is also coprime to $\ideal S=[\ideal S_1\cdots\ideal S_\lambda]$. Conversely, if $\ideal R$ and $\ideal S$ are coprime, then $\ideal R$ is coprime to each $\ideal S_j$. If $\ideal R=[\ideal R_1\cdots\ideal R_\mu]$ and every $\ideal R_i$ is coprime to every $\ideal S_j$, then $\ideal R$ and $\ideal S$ are coprime; here too the converse holds.}

Indeed, if $\ideal R$ is coprime to every $\ideal S_j$, then there are elements $s_j$ such that
\[
 s_j\equiv0(\ideal S_j),
 \qquad
 s_j\equiv\varepsilon(\ideal R).
\]
Consequently
\[
 s_1s_2\cdots s_\lambda\equiv0(\ideal S),
 \qquad
 s_1s_2\cdots s_\lambda\equiv\varepsilon(\ideal R),
 \qquad
 (\ideal R,
\ideal S)=(\varepsilon).
\]
Since $(\ideal R,
\ideal S)$ is divisible by each $(\ideal R,
\ideal S_j)$, the converse follows. Repeated application of the same argument gives the second part. For if $\ideal R_i$ is coprime to $\ideal S_1,
\ldots,\ideal S_\lambda$ for a fixed $i$, then $\ideal R_i$ is coprime to $\ideal S$; if this holds for every $i$, then, since coprimeness is mutual, $\ideal S$ is coprime to $\ideal R$. Conversely, coprimeness of $\ideal R$ and $\ideal S$ gives coprimeness of $\ideal S$ with $\ideal R_i$, and hence of $\ideal R_i$ with $\ideal S$.

The proof of existence and uniqueness\footnote{The existence of the decomposition can again be proved directly in analogy with \S~2; the uniqueness proof too can be carried out directly (compare what was said in the introduction about Schmeidler and Noether--Schmeidler). The proof given here also gives insight into the structure of coprime-irreducible ideals.} of the decomposition into coprime-irreducible ideals is, as in the corresponding proof for relatively-prime ideals, a matter of a unique grouping. Since, however, the relatively-prime irreducible ideals $\ideal R_1,
\ldots,\ideal R_\sigma$ of $\ideal M$ are uniquely defined by Theorem XII, it is unnecessary here to go back to the associated prime ideals.

We collect the uniquely defined relatively-prime irreducible ideals $\ideal R_1,
\ldots,\ideal R_\sigma$ of $\ideal M$ into groups so that \emph{every ideal of one group is coprime to every ideal of any different group, while a single group cannot be split into two subgroups such that every ideal of one subgroup is coprime to every ideal of the other subgroup}. Such a grouping is obtained as follows. By definition a group $G$ must contain, together with every ideal $\ideal R$, all ideals not coprime to $\ideal R$. Let these be $\ideal R_{i_1}$; let $\ideal R_{i_1i_2}$ be the ideals not coprime to these; in general let $\ideal R_{i_1\ldots i_\lambda}$ be not coprime to $\ideal R_{i_1\ldots i_{\lambda-1}}$. Since only finitely many ideals occur, this process terminates after finitely many steps. The set of ideals so obtained forms a group $G$ with the desired properties: all ideals outside $G$ are, by construction, coprime to all ideals in $G$. If $G$ were split into two subgroups $G^{(1)}$ and $G^{(2)}$, and if $\ideal R_{i_1\ldots i_\lambda}$ belonged to $G^{(1)}$, then, because coprimeness is mutual, $G^{(1)}$ would also have to contain $\ideal R_{i_1\ldots i_{\lambda-1}},
\ldots,\ideal R_{i_1}$, hence also $\ideal R$ and therefore the whole group $G$. This proves irreducibility. Repeating the construction with the remaining ideals gives a grouping $G_1,
\ldots,G_\tau$ of all $\ideal R$.

The grouping is unique: if $G'_1,
\ldots,G'_\tau$ is a second grouping and $\ideal R_{i_1\ldots i_\lambda}$ is an element of $G'_i$, then $G'_i$ contains $\ideal R$ and hence $G_1$, and, by the irreducibility condition, cannot contain elements different from $G_1$; therefore $G'_i=G_1$.

Let $\ideal T_i$ be the least common multiple of the ideals $\ideal R$ united in a group $G_i$. We show that
\[
 \ideal M=[\ideal T_1\cdots\ideal T_\tau]
\]
is a representation of $\ideal M$ by coprime-irreducible ideals. Decomposing the $\ideal T$ into the $\ideal R$ first shows that $[\ideal T_1\cdots\ideal T_\tau]$ really represents $\ideal M$. By Theorem XIV the $\ideal T$ are pairwise coprime, and each $\ideal T$ is coprime-irreducible. This proves the existence of such a representation.

For uniqueness, let $\ideal M=[\bar{\ideal T}_1\cdots\bar{\ideal T}_\tau]$ be a second representation of this type. Decomposing the $\bar{\ideal T}$ into relatively-prime irreducible ideals $\ideal R$, Theorem XIV shows that ideals $\ideal R$ occurring in different $\bar{\ideal T}$ are coprime to one another, hence also mutually relatively prime; therefore they agree with the uniquely defined relatively-prime irreducible ideals $\ideal R$ of $\ideal M$. The ideals $\bar{\ideal T}_i$ also induce, by Theorem XIV, a grouping $G'_i$ of the $\ideal R$ with the stated properties. Since this grouping is unique and each $\bar{\ideal T}_i$ is uniquely determined by the group $G'_i=G_i$, one has $\bar{\ideal T}_i=\ideal T_i$; uniqueness is proved.

For pairwise coprime ideals, the least common multiple is moreover equal to the product. For by Theorem XIV, if $\ideal M=[\ideal T_1\cdots\ideal T_\tau]$, the complement $\ideal L_i$ is also coprime to $\ideal T_i$; hence there are elements
\[
 t_i\equiv0(\ideal T_i),
 \qquad
 l_i\equiv0(\ideal L_i),
 \qquad
 \varepsilon=t_i+l_i.
\]
From $f\equiv0(\ideal T_i)$ and $f\equiv0(\ideal L_i)$ it follows that
\[
 f=f\varepsilon=ft_i+fl_i,
 \qquad
 f\equiv0(\ideal L_i\ideal T_i).
\]
Conversely $\ideal L_i\ideal T_i$ is divisible by $[\ideal L_i,
\ideal T_i]$, whence $[\ideal L_i,
\ideal T_i]=\ideal L_i\ideal T_i$; continuing the procedure on $\ideal L_i$ gives
\[
 \ideal M=[\ideal T_1\cdots\ideal T_\tau]
       =\ideal T_1\ideal T_2\cdots\ideal T_\tau.
\]
Therefore:

\textbf{Theorem XV.} \emph{Every ideal can be represented uniquely as a product of finitely many pairwise coprime and coprime-irreducible ideals.}

\providecommand{\eqzero}[1]{\equiv 0\,(#1)}
\providecommand{\ideal}[1]{\mathfrak{#1}}
\setcounter{footnote}{34}

\begin{center}
\S~9.\\[0.25em]
\textbf{Extension of the investigation to modules. Equality of the number of components in decompositions into irreducible modules.}
\end{center}

We now show that the content of the first three paragraphs, which concerns irreducible ideals and not primary and prime ideals, remains valid under weaker hypotheses. Those paragraphs do not use the commutative law of multiplication and refer only to the property of ideals of being modules. They therefore remain valid for modules with respect to non-commutative domains, which are now to be defined. The definition of modules is to be based on a double domain $(\Sigma,T)$ with the following properties.

\emph{$\Sigma$ is an abstractly defined non-commutative ring}; that is, $\Sigma$ is a system of elements $a,b,c,\ldots$ for which two modes of combination are defined, ring addition $(\#)$ and ring multiplication $(\times)$, satisfying the laws stated in \S~1, with the exception of the commutative law 4 for ring multiplication.

$T$ is a system of elements $\alpha,\beta,\gamma,\ldots$, for which, in connection with $\Sigma$, two combinations are likewise defined: addition, which from any two elements $\alpha,\beta$ produces a unique third element $\alpha+\beta$; and multiplication of an element $\alpha$ of $T$ by an element $c$ of $\Sigma$, which uniquely produces an element $c\alpha$ of $T$.\footnote{Thus here one is dealing with a ``right-sided'' multiplication, a ``right-sided'' domain $T$, and consequently with ``right-sided'' modules and ideals. If one were to base $T$ instead on a left-sided multiplication $\alpha c$, a corresponding theory of left-sided modules and ideals would result; $M$ would contain $\alpha c$ along with $\alpha$. The associative law would then have the form $(\gamma b)a=\gamma(b\times a)$.}

For these combinations the following laws hold:
\begin{enumerate}
\item The associative law of addition: $(\alpha+\beta)+\gamma=\alpha+(\beta+\gamma)$.
\item The commutative law of addition: $\alpha+\beta=\beta+\alpha$.
\item The law of unrestricted and unique subtraction: there is in $T$ one and only one element $\xi$ satisfying the equation $\alpha+\xi=\beta$ (one writes $\xi=\beta-\alpha$).
\item The associative law of multiplication: $a\cdot(b\cdot\gamma)=(a\times b)\cdot\gamma$.
\item The distributive law:
\[
 (a\# b)\cdot\gamma=a\cdot\gamma+b\cdot\gamma,\qquad c\cdot(\alpha+\beta)=c\cdot\alpha+c\cdot\beta.
\]
\end{enumerate}
From these conditions there follow, as is well known, the existence of the zero element and the validity of the distributive law also for the combination of subtraction and multiplication:
\[
 (a-b)\cdot\gamma=a\cdot\gamma-b\cdot\gamma,\qquad c\cdot(\alpha-\beta)=c\cdot\alpha-c\cdot\beta,
\]
where $(-)$ denotes subtraction in $\Sigma$. \emph{If $\Sigma$ contains a unit $\varepsilon$, then $\varepsilon\cdot\alpha=\alpha$ is to hold for every element $\alpha$ of $T$.}

By a module $M$ in $(\Sigma,T)$ we mean a system of elements of $T$ satisfying the two conditions:
\begin{enumerate}
\item \emph{Along with $\alpha$, $M$ contains $c\cdot\alpha$, where $c$ is any element of $\Sigma$.}
\item \emph{Along with $\alpha$ and $\beta$, $M$ contains the difference $\alpha-\beta$}; hence also $n\alpha$ for every integer $n$.\footnote{The integers are again to be understood as abbreviating symbols, not as ring elements.}
\end{enumerate}
By this definition, $T$ itself forms a module in $(\Sigma,T)$. If in particular the domain $T$ and the combinations established there coincide with the domain $\Sigma$ and the combinations valid there, then the module $M$ becomes a right ideal $M$ in $\Sigma$. If $\Sigma$ is also assumed commutative, the ordinary concept of ideal results, and is thus a special case of the concept of module.\footnote{The simplest example of a module is the module of integral linear forms: here $\Sigma$ consists of all rational integers and $T$ of all integral linear forms. A somewhat more general module arises when, in $\Sigma$ and $T$, one takes algebraic integers instead of rational integers, or for instance all even integers. If, instead of the linear forms, one considers each complex of coefficients as an element, then the combinations in $\Sigma$ and $T$ are in fact different. Ideals in non-commutative polynomial ring domains are the subject of the joint paper of Noether--Schmeidler. Ideals in other special non-commutative domains are treated in Hurwitz's lectures on the number theory of quaternions (Berlin, Springer 1919) and in the works of Du Pasquier cited there.}

For modules all definitions of \S~1 remain valid. Thus $\alpha\eqzero{M}$, respectively $N\eqzero{M}$, means that $\alpha$, respectively every element of $N$, is an element of $M$; in other words, $\alpha$, respectively $N$, is divisible by $M$. $M$ is a proper divisor of $N$ if $M$ contains elements different from those of $N$; from $N\eqzero{M}$ and $M\eqzero{N}$ follows $M=N$. The definitions of greatest common divisor and least common multiple likewise remain literally the same. If, in particular, the module $M$ contains a finite number of elements $\alpha_1,\ldots,\alpha_e$ such that $M=(\alpha_1\ldots\alpha_e)$, that is,
\[
 \alpha=c_1\alpha_1+\cdots+c_e\alpha_e+n_1\alpha_1+\cdots+n_e\alpha_e
\]
for every $\alpha\eqzero{M}$, where the $c_i$ are elements of $\Sigma$ and the $n_i$ are integers, then $M$ is called a \emph{finite module}, and $\alpha_1,\ldots,\alpha_e$ a \emph{module basis}.

In what follows we assume, just as in \S~1, only those domains $(\Sigma,T)$ that satisfy the finiteness condition: every module in $(\Sigma,T)$ is finite, and hence has a module basis.

For such a domain $(\Sigma,T)$, Theorem I on finite chains holds also for modules, as the proof given there shows, and therefore all hypotheses for \S\S~2 and 3 are satisfied. Definition I and Lemma I on shortest and reduced representations carry over directly; likewise Theorem II on the representability of every module as the least common multiple of finitely many irreducible modules, where Lemma II shows that every shortest representation of this kind is at the same time reduced. Theorem III also remains valid, expressing reducibility of a module through properties of its complement; from it follows Lemma III, and finally Theorem IV, which asserts equality of the number of components in two different shortest representations of a module as the least common multiple of irreducible modules. Theorem IV also gives Lemma IV as the converse of Lemma I on reduced representations.

\emph{Addendum.} The same argument shows that all these theorems and definitions remain valid if, for two-sided domains $T$, that is, domains which are both right- and left-sided, one everywhere understands by a module a two-sided module: one which, along with $\alpha$, also contains $c\alpha$ and $\alpha c$, and along with $\alpha$ and $\beta$ also contains $\alpha-\beta$.

While all these theorems rest only on the concept of divisibility and of least common multiple, the further uniqueness theorems rest essentially on the product concept and therefore do not admit a direct transfer. Thus the definition of a primary ideal and of a prime ideal cannot be transferred to modules, since the product of two elements of $T$ is not defined. The formally possible transfer to non-commutative rings loses its meaning, since here the existence of the associated prime ideal cannot be proved\footnote{For from $p_1^\rho\eqzero{\ideal Q}$ and $p_2^\sigma\eqzero{\ideal Q}$ it does not follow here that $(p_1-p_2)^{\rho+\sigma}\eqzero{\ideal Q}$; likewise, from $(ab)^\rho\eqzero{\ideal Q}$ it does not follow that $a^\rho b^\rho\eqzero{\ideal Q}$. Thus $\ideal P$ can be proved neither to be an ideal nor to have the property of a prime ideal. For two-sided ideals $\ideal P$ can indeed be proved to be an ideal, but not a prime ideal.} and the proof that an irreducible ideal is primary also fails. These two facts form the basis of the following uniqueness theorems. By contrast, if the non-commutative ring has a unit, coprime and coprime-irreducible ideals can be defined, and the arguments used in the proof of Theorem II show that every ideal can be represented as the least common multiple of finitely many pairwise coprime and coprime-irreducible ideals.\footnote{For special, ``completely reducible'' ideals, uniqueness theorems can be set up here too; compare the joint paper Noether--Schmeidler.}

Finally, we mention a sufficient criterion for the finiteness condition to hold in $(\Sigma,T)$: if $\Sigma$ contains a unit and satisfies the finiteness condition, and if $T$ itself is a finite module in $(\Sigma,T)$, then every module in $(\Sigma,T)$ is finite.

Indeed, from the existence of the unit in $\Sigma$, for
\[
 \ideal M=(f_1,
\ldots,f_e),
 \qquad
 f=\bar b_1f_1+\cdots+\bar b_ef_e+n_1f_1+\cdots+n_ef_e,
\]
where the $\bar b_i$ are elements of $\Sigma$ and the $n_i$ are integers, there follows also a representation
\[
 f=b_1f_1+\cdots+b_ef_e
\]
with the $b_i$ in $\Sigma$. For since $f_i=\varepsilon f_i$, one has $n_if_i=n_i\varepsilon f_i$, and since $n_i\varepsilon=(\varepsilon+\cdots+\varepsilon)$ belongs to $\Sigma$, $b_i=\bar b_i+n_i\varepsilon$ is an element of $\Sigma$. The hypothesis on $T$ says that every element $\alpha$ of $T$ admits a representation
\[
 \alpha=\bar a_1\tau_1+\cdots+\bar a_k\tau_k+n_1\tau_1+\cdots+n_k\tau_k,
\]
and hence
\[
 \alpha=a_1\tau_1+\cdots+a_k\tau_k,
\]
the second representation following from the first just as above because $\varepsilon\tau_i=\tau_i$.

As $\alpha$ runs through the elements of a module $M$ in $(\Sigma,T)$, the coefficients $a_k$ of $\tau_k$ run through an ideal $\ideal M_k$ of $\Sigma$; by the preceding, for every $a_k\eqzero{\ideal M_k}$ one has
\[
 a_k=b_1a_k^{(1)}+\cdots+b_ea_k^{(e)}.
\]
Let $\alpha^{(i)}$ be an element of $M$ for which the coefficient of $\tau_k$ is $a_k^{(i)}$. Then $\alpha-b_1\alpha^{(1)}-\cdots-b_e\alpha^{(e)}$ is an element belonging to $M$ which depends only on $\tau_1,
\ldots,\tau_{k-1}$. The procedure can be repeated for the totality of these elements, which no longer contain $\tau_k$ and which form a module $M'$, and the assertion is proved by finitely many repetitions.

\begin{center}
\S~10.\\[0.25em]
\textbf{Special case of the polynomial domain.}
\end{center}

1. Let the underlying ring domain $\Sigma$ consist of all polynomials in $x_1,
\ldots,x_n$ with arbitrary complex coefficients. By Hilbert's basis theorem (Ann. 36) the finiteness condition holds for it. The issue is the relation of our theorems to the known theorems of elimination theory and module theory.

This relation is established by the following special case of a known theorem of Hilbert:\footnote{Über die vollen Invariantensysteme. Math. Ann. 42 (1893), \S~3, p. 313.}

If $f$ vanishes for every finite system of values of $x_1,
\ldots,x_n$ that is a zero of all polynomials of a prime ideal $\ideal P$ (a zero of $\ideal P$), then $f$ is divisible by $\ideal P$. In other words, a prime ideal $\ideal P$ consists of the totality of polynomials which vanish at its zeros.\footnote{As Lasker showed [Math. Ann. 60 (1905), p. 607], this special case can also be proved directly for homogeneous forms; conversely, Hilbert's theorem follows again from it in the homogeneous and inhomogeneous cases. We can express that theorem as follows: if an ideal $\ideal R$ vanishes at all zeros of $M$, then a power of $\ideal R$ is divisible by $M$. This theorem, and also the special case, holds only if the range of values of the $x$ is algebraically closed; hence it cannot follow from our theorems alone, but must use the existence of roots, for example at the point where one uses that an ideal whose zeros consist only of $x_1=0,
\ldots,x_n=0$ contains all products of powers of the $x$ from some dimension on. The remaining proof can be somewhat simplified in comparison with Lasker by using our theorems. For Lasker must also appeal to Hilbert's theorem to prove the decomposition of an ideal into greatest primary ideals.}

Thus, if a product $fg$ vanishes for all zeros of $\ideal P$, at least one factor vanishes; the zeros form an irreducible algebraic variety. Conversely, if this definition of irreducible variety is adopted, then the totality of polynomials vanishing on an irreducible variety forms a prime ideal. Prime ideal and irreducible variety therefore correspond one-to-one. Since $\ideal Q\eqzero{\ideal P}$ and $\ideal P^e\eqzero{\ideal Q}$, the zeros of a primary ideal coincide with those of its associated prime ideal.\footnote{Macaulay (compare the introduction) takes this property of a primary ideal, namely to possess an irreducible variety, as the definition; Lasker includes only the concept of the manifold of a variety in the definition and otherwise defines abstractly. The primary ideals vanishing only for $x_1=0,
\ldots,x_n=0$ occupy a special position for Lasker.} Therefore the representation of an ideal as the least common multiple of greatest primary ideals gives a resolution of all zeros of the ideal into irreducible varieties; and, as Lasker showed, the converse also holds. The proof of uniqueness of the associated prime ideals thus corresponds here to the fundamental theorem of elimination theory on the unique decomposability of an algebraic variety into irreducible varieties, and for special polynomial domains, where no unique product representation of the polynomials by irreducible polynomials of the domain and consequently no elimination theory exists, it can serve as an equivalent of that theorem of elimination theory.

The irreducible varieties corresponding to the isolated primary ideals are exactly those which occur in the ``minimal resolvent'';\footnote{Compare, for example, J. König, Einleitung in die allgemeine Theorie der algebraischen Größen (Leipzig, Teubner, 1903), p. 235.} for the zeros of every non-isolated greatest primary ideal are at the same time zeros of at least one isolated one, namely one whose associated prime ideal is divisible by that of the non-isolated ideal. The uniqueness of isolated primary ideals therefore gives, in the exponents, new invariant multiplicity numbers. The uniqueness theorems for the resolution of primary ideals into irreducible ones can likewise be regarded as supplements to elimination theory in the sense of multiplicity.

After these remarks the different decompositions can be interpreted in terms of their behavior with respect to algebraic varieties. Pairwise coprime ideals correspond to varieties having no common zero; for mutually relatively prime ideals, no irreducible variety of one ideal is at the same time a common zero; the greatest primary ideals vanish only on irreducible varieties which are all distinct from one another; in the decomposition into irreducible ideals the same irreducible varieties may occur with multiplicity.

It should also be noted that, instead of the general polynomial domain, the domain of all homogeneous forms may also be taken as the basis; for one readily verifies that the general theorems remain valid for the combinations there in force, where addition is defined only for forms of the same dimension.\footnote{That in the ambiguous case inhomogeneous decompositions can exist beside homogeneous ones is shown by the example $(x^3,xy,y^3)=[(x^3,y),(y^3,x)]=[(xy,x^3,y^3,x+y^3),(xy,x^3,y^3,y+x^2)]$.}

A simplest example of the four different decompositions - for which the formulas follow below - is given, according to the preceding, by a line and two intersecting lines skew to it, one of the latter containing, with higher multiplicity, a point different from the point of intersection. The decomposition into coprime-irreducible ideals corresponds to the decomposition into the line and the variety skew to it; in the decomposition into relatively-prime irreducible ideals this variety splits into the two lines; the decomposition into greatest primary ideals corresponds to separating off the point of higher multiplicity, while the decomposition into irreducible ideals forces a resolution of that point.

Taking this point as the origin, the running line as the $y$-axis, the line intersecting it parallel to the $x$-axis, and the line skew to it parallel to the $z$-axis, such a configuration is represented, for example, by the following \emph{irreducible ideals}\footnote{Of these, the first three are irreducible as prime ideals; $\ideal B_4$ because it has only the divisors $(x^2,y,z)$ and $(x,y,z)$; $\ideal B_5$ because every divisor contains the polynomial $xy$.}:
\par\medskip
\noindent\makebox[\linewidth][c]{\(\displaystyle
 \ideal B_1=(x-1,y);\qquad
 \ideal B_2=(y-1,z);\qquad
 \ideal B_3=(x,z);\qquad
 \ideal B_4=(x^3,y,z);
\)}\par
\noindent\makebox[\linewidth][c]{\(\displaystyle
 \ideal B_5=(x^2,y^2,z).
\)}\par\medskip
The \emph{associated prime ideals} are:
\par\smallskip
\noindent\makebox[\linewidth][c]{\(\displaystyle
 \ideal P_1=\ideal B_1;\qquad
 \ideal P_2=\ideal B_2;\qquad
 \ideal P_3=\ideal B_3;\qquad
 \ideal P_4=\ideal P_5=(x,y,z).
\)}\par\medskip
The \emph{greatest primary ideals} are:
\par\smallskip
\noindent\makebox[\linewidth][c]{\(\displaystyle
 \ideal Q_1=\ideal B_1;\qquad
 \ideal Q_2=\ideal B_2;\qquad
 \ideal Q_3=\ideal B_3;\qquad
 \ideal Q_4=[\ideal B_4,\ideal B_5]=(x^3,y^2,x^2y,z).
\)}\par\medskip
The \emph{relatively-prime irreducible ideals} are:
\par\smallskip
\noindent\makebox[\linewidth][c]{\(\displaystyle
 \ideal R_1=\ideal Q_1;\qquad
 \ideal R_2=\ideal Q_2;\qquad
 \ideal R_3=[\ideal Q_3,\ideal Q_4]=(x^3,x^2y,xy^2,z).
\)}\par\medskip
The \emph{coprime-irreducible ideals} are:
\par\smallskip
\noindent\makebox[\linewidth][c]{\(\displaystyle
 \ideal S_1=\ideal R_1;\qquad
 \ideal S_2=[\ideal R_2,\ideal R_3]=((y-1)x^3,(y-1)x^2y,(y-1)xy^2,z).
\)}\par\medskip
This gives the \emph{total ideal}:
\par\smallskip
\noindent\(\displaystyle \ideal M=[\ideal S_1,\ideal S_2]=\ideal S_1\cdot\ideal S_2\)\par
\noindent\(\displaystyle
 =((x-1)(y-1)x^3,(y-1)x^2y,(y-1)xy^2,(x-1)z,y(y-1)x^3,yz),
\)\par\smallskip
\noindent\begin{tabular*}{\linewidth}{@{}l@{\extracolsep{\fill}}c@{}}
where & \(\displaystyle 1=-(y-1)x^3+(y-1)(x^3-1)+y,\)
\end{tabular*}\par
\noindent\makebox[\linewidth][c]{\(\displaystyle
 (y-1)(x^3-1)+y\eqzero{\ideal S_1};\qquad -(y-1)x^3\eqzero{\ideal S_2}.
\)}\par
Here the ideals $\ideal B_1$, $\ideal B_2$, and $\ideal B_3$ are isolated, and hence uniquely determined also in the decompositions into irreducible and greatest primary ideals. The ideals $\ideal B_4$ and $\ideal B_5$, respectively $\ideal Q_4$, are non-isolated; they are not uniquely determined, but can, for example, be replaced by
\par\smallskip
\noindent\makebox[\linewidth][c]{\(\displaystyle
 \ideal D_4=(x^3,y+\lambda x^2,z);\qquad
 \ideal D_5=(x^2+\mu xy,y^2,z).
\)}\par
Likewise, $\ideal Q_4$ can be replaced by
\par\smallskip
\noindent\makebox[\linewidth][c]{\(\displaystyle
 \overline{\ideal Q}_4=(x^3,x^2y,y^2+\lambda xy,z).
\)}\par

2. Just as the general, and the integral, polynomial domains satisfy the finiteness condition, so does every finite integral domain of polynomials, by Hilbert's theorem on module bases;\footnote{Conversely, if a polynomial domain satisfies the finiteness condition and every polynomial has at least one representation in which the multipliers are of lower degree in $x$, then the domain is a finite integral domain.} the coefficients may be assigned to an arbitrary field. We give one more example, for the domain of all even polynomials, the simplest domain in which, because $x^2y^2=(xy)^2$, no unique product representation of the polynomials by irreducible polynomials of the domain exists. It is the same configuration as in the preceding example, now given by the irreducible ideals
\[
\begin{gathered}
 \ideal B_1=(x^2-1,xy,y^2,yz),\quad
 \ideal B_2=(y^2-1,xz,yz,z^2),\\
 \ideal B_3=(x^2,xy,xz,yz,z^2),\quad
 \ideal B_4=(x^4,x^3y,y^2,xz,yz,z^2),\\
 \ideal B_5=(x^2,y^2,xz,yz,z^2).
\end{gathered}
\]
The associated prime ideals are
\[
 \ideal P_1=\ideal B_1,
 \quad \ideal P_2=\ideal B_2,
 \quad \ideal P_3=\ideal B_3,
 \quad \ideal P_4=\ideal P_5=(x^2,xy,y^2,xz,yz,z^2).
\]
That $\ideal P_1$ is a prime ideal follows from the fact that every polynomial of the domain has the form
\[
 f=\varphi(z^2)+xz\,\psi(z^2)\pmod{\ideal P_1}.
\]
Thus if
\[
 f_1=\varphi_1(z^2)+xz\,\psi_1(z^2),
 \qquad
 f_2=\varphi_2(z^2)+xz\,\psi_2(z^2),
\]
then from $f_1f_2\eqzero{\ideal P_1}$ one obtains the equations
\[
 \varphi_1\varphi_2+z^2\psi_1\psi_2=0,
 \qquad
 \varphi_2\psi_1+\varphi_1\psi_2=0,
\]
and hence $f_1\eqzero{\ideal P_1}$ or $f_2\eqzero{\ideal P_1}$. In exactly the same way one sees that $\ideal P_2$ is a prime ideal; $\ideal P_3$ is a prime ideal because every polynomial of the domain modulo $\ideal P_3$ is congruent to a polynomial in $y^2$; $\ideal P_4$ consists of all polynomials of the domain. Thus $\ideal B_1$, $\ideal B_2$, and $\ideal B_3$ are irreducible as prime ideals, while $\ideal B_4$ and $\ideal B_5$ each have only the single proper divisor $\ideal P_4$, and hence are necessarily irreducible.

From the irreducible ideals one obtains the greatest primary ideals
\[
\begin{gathered}
 \ideal Q_1=\ideal B_1,
 \quad \ideal Q_2=\ideal B_2,
 \quad \ideal Q_3=\ideal B_3,
 \quad \ideal Q_4=[\ideal B_4,\ideal B_5]=(x^4,x^3y,y^2,xz,yz,z^2),
\end{gathered}
\]
the relatively-prime irreducible ideals
\[
 \ideal R_1=\ideal Q_1,
 \quad \ideal R_2=\ideal Q_2,
 \quad
 \ideal R_3=[\ideal Q_3,\ideal Q_4]=(x^4,x^3y,x^2y^2,xy^3,xz,yz,z^2),
\]
and the coprime-irreducible ideals
\[
 \ideal S_1=\ideal R_1,
 \quad \ideal S_2=[\ideal R_2,\ideal R_3]
 =((y^2-1)x^4,(y^2-1)x^3y,(y^2-1)x^2y^2,
 (y^2-1)xy^3,xz,yz,z^2).
\]
As in the first example,
\[
 [\ideal B_4,\ideal B_5]=[\ideal D_4,\ideal D_5],
 \qquad
 [\ideal Q_3,\ideal Q_4]=[\ideal Q_3,
\overline{\ideal Q}_4],
\]
where
\[
 \ideal D_4=(x^4,xy+\lambda x^2,\ldots),
 \qquad
 \ideal D_5=(x^2+\mu xy,\ldots),
 \qquad \lambda\mu\ne1,
\]
and
\[
 \overline{\ideal Q}_4=(x^4,x^3y,y^2+\lambda xy,\ldots).
\]
The remaining irreducible, respectively greatest primary, ideals $\ideal B_1$, $\ideal B_2$, $\ideal B_3$ are uniquely determined as isolated ideals.

\begin{center}
\S~11.\\[0.25em]
\textbf{Examples from Number Theory and the Theory of Differential Expressions.}
\end{center}

1. Let the underlying domain $\Sigma$ consist of all even rational integers. Thus $\Sigma$ can be put in one-to-one correspondence with the domain of all rational integers, by assigning to each number $2a$ in $\Sigma$ the number $a$. It follows at once that every ideal in $\Sigma$ is a principal ideal $(2a)$, where in the basis representation $2c=n\cdot2a$ for every element $2c$ of the ideal, the odd integers $n$ are only abbreviations for finite sums.

The prime ideals of the domain are $\ideal P_0=\ideal D=(2)$ and $\ideal P=(2p)$, where $p$ denotes an odd prime number; every prime ideal is therefore divisible by $\ideal P_0$, but by no further prime ideal. The primary ideals are
\[
 \ideal Q_{e_0}=(2\cdot2^{e_0})
 \qquad\text{and}\qquad
 \ideal Q_e=(2p^e).
\]
They are at the same time irreducible ideals, and by what was said about the prime ideals, two ideals belonging to different odd prime numbers are mutually relatively prime, but no $\ideal Q$ is relatively prime to any $\ideal Q_{e_0}$.\footnote{Indeed, from $2b\cdot2p_1^{e_1}\eqzero{2p_2^{e_2}}$ for odd $p_1\ne p_2$ one always obtains $2b\eqzero{2p_2^{e_2}}$; but from $2b\cdot2p^e\eqzero{2\cdot2^{e_0}}$ one obtains only $2b\equiv2\cdot2^{e_0-1}\,(2\cdot2^{e_0})$.} Thus the $\ideal Q_{e_0}$ are the only non-isolated primary ideals. The unique decomposition of $a$ into prime powers corresponds to the unique representation of the ideal $(2a)$ by greatest primary, at the same time irreducible, ideals:
\[
 (2a)=[(2\cdot2^{e_0}),(2p_1)^{e_1},\ldots,(2p_\alpha)^{e_\alpha}].
\]
Thus here, in contrast to the examples from the polynomial domain, even the non-isolated greatest primary ideals are uniquely determined. If $e_0=0$ this is at the same time a representation by mutually relatively prime ideals, while if $e_0>0$ the ideal is relatively-prime-irreducible.

Thus, while in the domain of all integers the four different decompositions coincide, here this happens only for the two decompositions into greatest primary and into irreducible ideals on the one hand, and, for $e_0>0$, for coprime-irreducible (every ideal is coprime-irreducible, since the domain has no unit) and relatively-prime-irreducible ideals on the other; for $e_0=0$, however, the coprime-irreducible and relatively-prime-irreducible decompositions become different. At the same time one already obtains here an example in which one prime ideal can be divisible by another without being identical with it; more generally, one sees that product representation does not follow from divisibility. The latter fact - a consequence of the absence of a unit in the domain - is also why no unique product representation of the numbers of the domain by irreducible numbers of the domain exists, even though every ideal is a principal ideal; the introduction of least common multiples therefore proves necessary. It should also be noted that the situation remains exactly the same if, instead of all even integers, one takes all integers divisible by a fixed prime number or power of a prime.

Irreducible and primary ideals also become different, however, if $\Sigma$ consists of all numbers divisible by a composite number $g=p_1^{\sigma_1}\cdots p_\nu^{\sigma_\nu}$. Again every ideal is a principal ideal $(g\cdot a)$, and the prime ideals are again $\ideal P_0=\ideal D=(g)$ and $\ideal P=(g\cdot p)$, where $p$ is a prime different from the primes occurring in $g$. But the irreducible ideals are $\ideal B_{\lambda_i}=(g\cdot p_i^{\lambda_i})$ and $\ideal B_e=(g\cdot p^e)$; the primary ideals are $\ideal Q_e=\ideal B_e$ and the ideals, different from the irreducible ones,
\[
 \ideal Q_{\lambda_1\ldots\lambda_\nu}=(g\cdot p_1^{\lambda_1}\cdots p_\nu^{\lambda_\nu}),
\]
where the $\ideal B_{\lambda_i}$ and $\ideal Q_{\lambda_1\ldots\lambda_\nu}$ all have the same associated prime ideal $\ideal P_0=(g)$. Here too the decomposition into irreducible ideals is unique, and consequently so is the decomposition into greatest primary ideals; thus here again the non-isolated ideals are uniquely determined.

2. An example of a non-commutative ring domain is supplied by the ideal theory in non-commutative polynomial domains treated in the paper of Noether--Schmeidler. In particular one has ``completely reducible'' ideals, that is, ideals for which the components of the decomposition are pairwise coprime and have no proper divisors; the components are therefore a fortiori irreducible. Thus, in addition to the isomorphism proved there according to \S~9, one obtains equality of the number of components in two different decompositions. For the decomposition of systems of partial or ordinary linear differential expressions, which arises as a special case of that paper, this gives a result which does not seem to have been noticed even in the known case of a single ordinary linear differential expression.

At the same time the system $T$ of all residue classes of a fixed ideal $\ideal M$, together with the non-commutative polynomial domain $\Sigma$, supplies a double domain $(\Sigma,T)$, where $T$ has the module property with respect to $\Sigma$. For the difference of two residue classes is again a residue class, and likewise the product of a residue class by an arbitrary polynomial; whereas the product of two residue classes does not exist (loc. cit., \S~3). Thus the systems of residue classes called ``subgroups'' there furnish examples of modules in double domains $(\Sigma,T)$, where the underlying ring domain $\Sigma$ is non-commutative.

\begin{center}
\S~12.\\[0.25em]
\textbf{Example from Elementary-Divisor Theory.}
\end{center}

This is a conception of elementary divisor theory dictated by the general developments, but the theory itself is assumed as known.

Let $\Sigma$ be the domain of all integral matrices with $n^2$ elements, with addition and multiplication defined in the usual sense for matrices. Then $\Sigma$ is a non-commutative ring domain; the ideals are therefore in general one-sided, with two-sided ideals only as a special case.\footnote{The ideal theory of these domains is the subject of the works of Du Pasquier: Zahlentheorie der Tettarionen, dissertation, Zürich, Vierteljahrsschr. d. Naturf. Ges. Zürich, 51 (1906); Zur Theorie der Tettarionenideale, ibid., 52 (1907). The content of the second paper is the proof that every ideal is a principal ideal.}

We first show that every ideal is principal. For this purpose, in the case of right ideals, to each matrix $A=(a_{ik})$ we assign the module
\[
 A=(a_{11}\xi_1+\cdots+a_{1n}\xi_n,
 \ldots,
 a_{n1}\xi_1+\cdots+a_{nn}\xi_n)
\]
of integral linear forms. Conversely, to this module there corresponds every matrix that supplies a basis of $A$, hence, beside $A$, also $UA$, where $U$ is unimodular. More generally, to the product $PA$ there corresponds a module $B$ which is a multiple of $A$. A single linear form from $A$ is given by such matrices $P$ as contain only one non-zero row.

Let now $A_1,A_2,
\ldots,A_\nu,
\ldots$ be all elements of an ideal $\ideal M$; let $A_1,A_2,
\ldots,A_\nu,
\ldots$ be the modules assigned to them; let $A$ be their greatest common divisor and $UA$ the most general matrix assigned to this module $A$. To each single linear form in $A$ there then corresponds, by the definition of greatest common divisor, a matrix $P_1A_{i_1}+\cdots+P_\sigma A_{i_\sigma}$, where, as above, the $P$ have only one non-zero row. It follows that the matrix $A$ corresponding to a basis of $A$ also admits such a representation, now with general $P$, and hence becomes an element of $\ideal M$. Since, furthermore, every module $A_i$ is divisible by $A$, every matrix $A_i$ is divisible by $A$; $A$, and in general $UA$, forms a basis of $\ideal M$. If one is dealing with left ideals, then correspondingly the columns of each matrix are to be regarded as the basis of a module; every ideal is a principal ideal, for which, beside $A$, also $AV$ is a basis, with $V$ an arbitrary unimodular matrix.

For the connection with elementary divisor theory, we now base the following on two-sided ideals; hence in particular $PAQ$ belongs to the ideal whenever $A$ does. By the preceding, the most general basis of such an ideal is $UAV$, where $U$ and $V$ are unimodular. The basis elements therefore exhaust a class of equivalent matrices,\footnote{The basis elements of one-sided ideals correspond to right or left classes.} and there is a one-to-one relation between ideal and class, hence also between ideal and the elementary divisor system $(a_1\mid a_2\mid\cdots\mid a_n)$ of the class, where the $a_i$ are, as is well known, non-negative integers, each dividing the following one. Thus the matrix of the class that occurs in the normal form determined by the elementary divisors may be regarded as a special basis of the ideal; divisibility of ideals, respectively of classes, implies divisibility of the elementary divisors, and conversely.

But Du Pasquier has shown, loc. cit. \S~11, that for every two-sided ideal the rank is $n$ and all elementary divisors agree. To be able to consider the case of general elementary divisors, we must therefore start not from the ideals but directly from the two-sided classes (which are likewise to be denoted by capital German letters). A class $\mathfrak A=UAV$ is divisible by another class $\mathfrak B$ if $A=PBQ$.

In general: \emph{the least common multiple (the greatest common divisor) of two classes is obtained by forming the least common multiple (the greatest common divisor) of the corresponding elementary divisor systems}.\footnote{In the paper Zur Theorie der Moduln, Math. Ann. 52 (1899), p. 1, E. Steinitz defines the least common multiple (the greatest common divisor) of classes by the least common multiples and greatest common divisors of the elementary systems. Independently, the least common multiple of classes occurs as ``congruence composition'' in H. Brandt, Komposition der binären quadratischen Formen relativ einer Grundform, J. f. M. 150 (1919), p. 1.} For let
\[
 (a_1\mid a_2\mid\cdots\mid a_n),\quad
 (b_1\mid b_2\mid\cdots\mid b_n),\quad
 (c_1\mid c_2\mid\cdots\mid c_n),
\]
where $c_i=[a_i,b_i]$, be respectively the elementary divisor systems of $\mathfrak A$, $\mathfrak B$, and $\mathfrak C^*$, and let $\mathfrak C=[\mathfrak A,\mathfrak B]$. Then $\mathfrak C^*$ is divisible by $\mathfrak A$ and $\mathfrak B$, hence by $\mathfrak C$; conversely, the elementary divisors of $\mathfrak C$ are divisible by those of $\mathfrak C^*$, so $\mathfrak C$ is divisible by $\mathfrak C^*$ and therefore $\mathfrak C=\mathfrak C^*$. The proof for the greatest common divisor is obtained correspondingly.

The unique representation of the elementary divisors $a_i$ as least common multiples of prime powers therefore corresponds to a representation of $\mathfrak A$ as the least common multiple of classes $\mathfrak Q$ whose elementary divisors are given by powers of one prime number, in symbols
\[
 \mathfrak Q\sim(p^{r_1}\mid p^{r_2}\mid\cdots\mid p^{r_e}\mid0\mid\cdots\mid0),
 \qquad r_1\leqq r_2\leqq\cdots\leqq r_e.
\]
If in particular the rank $\varrho$ is equal to $n$, then one has here a decomposition into coprime and coprime-irreducible classes which is unique despite the non-commutative domain.

The classes $\mathfrak Q$ can be further decomposed into classes corresponding to the elementary divisor systems
\[
\begin{gathered}
 \mathfrak B_1\sim(p^{r_1}\mid\cdots\mid p^{r_1}),\quad
 \mathfrak B_2\sim(1\mid p^{r_2}\mid\cdots\mid p^{r_2}),\quad \ldots,\\
 \mathfrak B_e\sim(1\mid\cdots\mid1\mid p^{r_e}\mid\cdots\mid p^{r_e}),\quad
 \mathfrak B_{e+1}\sim(1\mid\cdots\mid1\mid0\mid\cdots\mid0),
\end{gathered}
\]
where in $\mathfrak B_\nu$ the number $1$ occurs $(\nu-1)$ times. If, for example, $r_1=r_2=\cdots=r_\mu=0$, then $\mathfrak B_1,
\ldots,\mathfrak B_\mu$ are equal to the unit class and are to be omitted from the decomposition; the same holds for $\mathfrak B_{e+1}$ if $\varrho=n$. If, furthermore, say $r_\nu=r_{\nu+1}=\cdots=r_{\nu+\lambda}$, then $\mathfrak B_{\nu+1},
\ldots,\mathfrak B_{\nu+\lambda}$ are proper divisors of $\mathfrak B_\nu$, and are likewise to be discarded. Denote by $\mathfrak B_{i_1},
\ldots,\mathfrak B_{i_k}$ those that remain and now give a shortest representation. Then
\[
 \mathfrak Q=[\mathfrak B_{i_1},
\ldots,\mathfrak B_{i_k}]
\]
is the unique decomposition of $\mathfrak Q$ into irreducible classes. Indeed, suppose
\[
 \mathfrak B_\nu\sim(1\mid\cdots\mid1\mid p^{r_\nu}\mid\cdots\mid p^{r_\nu})
\]
can be represented as the least common multiple of
\[
 \mathfrak C\sim(1\mid\cdots\mid1\mid p^{s_1}\mid\cdots\mid p^{s_\lambda})
 \quad\text{and}\quad
 \mathfrak D\sim(1\mid\cdots\mid1\mid p^{t_1}\mid\cdots\mid p^{t_\mu}).
\]
Then in the elementary divisor systems of $\mathfrak C$ and $\mathfrak D$ the number $1$ must stand in the first $(\nu-1)$ places; in the $\nu$-th place one of the exponents $s_\nu$ or $t_\nu$, say $s_\nu$, must become equal to $r_\nu$. But since
\[
 r_\nu=s_1\leqq s_2\leqq\cdots\leqq s_\lambda\leqq r_\nu,
\]
one obtains $\mathfrak C=\mathfrak B_\nu$, and $\mathfrak B_\nu$ is therefore irreducible.\footnote{The $\mathfrak B$ are still reducible into one-sided classes; here there is no longer a unique relation between elementary divisors and class, and hence no unique decomposition into irreducible one-sided classes. The following example, supplied to me by H. Brandt, shows this for decomposition into right-sided classes (where the classes are represented by a basis matrix, or by the corresponding module):
\[
\begin{gathered}
 \mathfrak B=[\mathfrak C_1,\mathfrak C_2]=[\mathfrak D_1,\mathfrak D_2],\quad
 \mathfrak B\sim\begin{pmatrix}p&0\\0&p\end{pmatrix},\quad
 \mathfrak C_1\sim\begin{pmatrix}1&0\\0&p\end{pmatrix},\quad
 \mathfrak C_2\sim\begin{pmatrix}p&0\\0&1\end{pmatrix},\\
 \mathfrak D_1\sim\begin{pmatrix}1&0\\0&p\end{pmatrix},
 \quad
 \mathfrak D_2\sim\begin{pmatrix}p&0\\(p-1)&1\end{pmatrix}.
\end{gathered}
\]
Indeed, the modules $(\xi,p\eta)$, $(p\xi,\eta)$, and $(p\xi,(p-1)\xi+\eta)$ each have as their only proper divisor the module $(\xi,\eta)$, and are therefore irreducible and mutually distinct.} The same holds for $\mathfrak B_{e+1}$, where $p$ is to be replaced by $0$. But, as the formation of the least common multiple shows, each of these irreducible classes gives a definite exponent and the place where this exponent first occurs in the elementary divisor system of $\mathfrak Q$, respectively $\mathfrak A$, while $\mathfrak B_{e+1}$ gives the rank. Since these numbers are uniquely determined by the elementary divisor system of $\mathfrak Q$, respectively $\mathfrak A$, and since the relation between elementary divisors and class is one-to-one, the decomposition of $\mathfrak Q$, and likewise of any class, into irreducible classes is unique.

In summary: \emph{Every two-sided class $\mathfrak A$ of integral matrices with bounded number of entries can be represented uniquely as the least common multiple of finitely many irreducible two-sided classes. Each irreducible class represents a fixed prime divisor of the elementary divisor system of $\mathfrak A$, an associated exponent, and the place where that exponent first occurs. The irreducible class corresponding to the divisor $0$ gives the rank of $\mathfrak A$.}

\begin{flushleft}
Erlangen, October 1920.
\end{flushleft}
\begin{center}
(Received 16. 10. 1920.)
\end{center}


\end{document}
